% ---------------------------------------------------------------------------
%  bdsim -- mathematical and algorithmic reference
%  Compile with:  pdflatex theory.tex   (twice, for cross-references)
% ---------------------------------------------------------------------------
\documentclass[11pt,a4paper]{report}

\usepackage[margin=2.5cm]{geometry}
\usepackage{amsmath,amssymb,bm}
\usepackage{booktabs}
\usepackage{longtable}
\usepackage{array}
\usepackage{listings}
\usepackage{xcolor}
\usepackage[hidelinks]{hyperref}

\lstset{
  basicstyle=\ttfamily\small,
  breaklines=true,
  frame=single,
  framesep=4pt,
  rulecolor=\color{gray!50},
  backgroundcolor=\color{gray!5},
  columns=fullflexible,
}

\newcommand{\D}{\bm{D}}
\newcommand{\B}{\bm{B}}
\newcommand{\R}{\bm{R}}
\newcommand{\Q}{\bm{Q}}
\newcommand{\uv}{\bm{u}}
\newcommand{\dt}{\Delta t^{*}}
\newcommand{\kB}{k_\mathrm{B}}
% Long paths/identifiers in \code should be breakable, otherwise they overrun
% the margin. Insert discretionary breaks after / and _ inside typewriter text.
\newcommand{\code}[1]{\texttt{\seqsplitcode{#1}}}
\newcommand{\seqsplitcode}[1]{#1}
\setlength{\emergencystretch}{3em}
\sloppy

\title{\textbf{bdsim}\\[0.4em]
       \large Mathematical and algorithmic reference for the\\
       C\texttt{++}/Python Brownian dynamics bead--spring code}
\author{}
\date{\today}

\begin{document}
\maketitle

\begin{abstract}
\noindent
This document specifies the mathematics implemented by the \code{bdsim} code: the
stochastic differential equation that is integrated, the semi-implicit
predictor--corrector scheme, every spring force law and how its implicit
per-step equation is solved, the hydrodynamic-interaction machinery (RPY tensor,
Cholesky and Chebyshev square roots, spectral bounds), the bending and
excluded-volume potentials, the random number generator, the initial-configuration
samplers, user-applied external forces, and the measured observables. It also
covers the inverse problem: how a real semiflexible molecule of given contour and
persistence length is coarse-grained onto the model, which moments of the
wormlike-chain distribution that fit should target, and what the resulting chain
does and does not reproduce --- in particular its force-extension behaviour. The
dynamics are matched separately, by fixing the bead hydrodynamic radius from a
measured diffusivity, which also supplies the conversions between simulation and
laboratory units. A final chapter covers the estimation of error bars from
correlated time series, without which none of the comparisons mean anything.

It is derived from the appendices of the author's thesis, but has been edited to
describe \emph{the code as it currently stands}. Where the implementation departs
from the thesis --- most importantly in the estimation of the spectral bounds used
by the Chebyshev square root, and in several robustness measures added to the
implicit spring solve --- the departure is stated explicitly and the reason given.
Section~\ref{sec:departures} collects these differences in one place.
\end{abstract}

\tableofcontents

% ===========================================================================
\chapter{Scope, notation and units}
\label{ch:scope}

\section{What the code does}

\code{bdsim} integrates the Brownian dynamics of a single bead--spring chain of
$N$ beads connected by $N-1$ springs, optionally with hydrodynamic interaction
(HI), excluded volume (EV) and a bending potential, in a homogeneous flow field.
The core is C\texttt{++}; a Python layer (via \code{nanobind}) owns run
configuration, initial configurations, sampling, ensemble averaging and
input/output.

\section{Non-dimensionalisation}
\label{sec:units}

All quantities are in the standard Hookean non-dimensional form. With spring
constant $H$, bead friction coefficient $\zeta$ and thermal energy $\kB T$, the
scales are
\begin{equation}
  l_H = \sqrt{\frac{\kB T}{H}}, \qquad
  \lambda_H = \frac{\zeta}{4H}, \qquad
  F_H = \sqrt{H \kB T},
\end{equation}
for length, time and force respectively. Starred quantities are non-dimensional;
the star is dropped wherever it is unambiguous. In these units $H = 1$ and
$\kB T = 1$.

The hydrodynamic interaction parameter is
\begin{equation}
  h^{*} = \frac{a}{\sqrt{\pi}\, l_H},
  \qquad\text{equivalently}\qquad
  a^{*} = \sqrt{\pi}\, h^{*},
  \label{eq:hstar}
\end{equation}
where $a$ is the bead hydrodynamic radius. Equation~\eqref{eq:hstar} is worth
stating explicitly because it fixes the crossover radius in the RPY tensor
(Section~\ref{sec:rpy}): the near-field branch is entered when the bead
separation falls below $2a^{*} = 2\sqrt{\pi}\,h^{*}$.

\section{Notation}

\begin{longtable}{@{}p{0.30\textwidth}p{0.64\textwidth}@{}}
\toprule
Symbol & Meaning \\
\midrule
\endhead
$N$ & number of beads \\
$\bm{r}_\nu$ & position of bead $\nu$, $\nu = 1,\dots,N$ \\
$\R$ & the $3N$ block vector of all bead positions \\
$\Q_\mu = \bm{r}_{\mu+1}-\bm{r}_\mu$ & connector (bond) vector $\mu$, $\mu = 1,\dots,N-1$ \\
$Q_\mu = \|\Q_\mu\|$ & bond length \\
$\uv_\mu = \Q_\mu/Q_\mu$ & bond unit vector \\
$\theta_\mu$ & angle between bonds $\mu-1$ and $\mu$ \\
$\bm{K}$ & velocity-gradient tensor $\bm{\kappa}$, applied blockwise \\
$\D$ & $3N\times3N$ diffusion (mobility) tensor \\
$\B$ & a matrix with $\D = \B\cdot\B^{\mathsf{T}}$ \\
$\bm{F}^{\mathrm{S}}$ & spring force on each bead (treated implicitly) \\
$\bm{F}^{\mathrm{E}}$ & all non-spring forces: EV $+$ bending (explicit) \\
$\bm{F}^{\mathrm{c}}_\mu$ & force carried by spring $\mu$ \\
$\dt$ & timestep \\
$n$, $j$ & timestep index, implicit-iteration index \\
\bottomrule
\end{longtable}

A tilde denotes a predicted (Euler) quantity, e.g.\ $\tilde{\R}_{n+1}$.

\section{Source map}

\begin{longtable}{@{}p{0.34\textwidth}p{0.60\textwidth}@{}}
\toprule
File & Contents \\
\midrule
\endhead
\code{src/integrator.cpp} & the predictor--corrector scheme (Chapter~\ref{ch:integrator}) \\
\code{src/spring.cpp}     & force laws and the implicit connector solve (Chapter~\ref{ch:springs}) \\
\code{src/numerics.cpp}   & cubic root finder, Newton polish, \code{rtsafe} \\
\code{src/hydrodynamics.cpp} & RPY tensor, Chebyshev square root, spectral bounds (Chapter~\ref{ch:hi}) \\
\code{src/mobility.cpp}   & \code{Diffusion}/\code{Mobility}, Brownian displacement, noise \\
\code{src/linalg.cpp}     & dense mat--vec, Cholesky factorisation \\
\code{src/bending.cpp}    & bending force (Chapter~\ref{ch:bending}) \\
\code{src/excluded\_volume.cpp} & EV force (Chapter~\ref{ch:ev}) \\
\code{src/flow.cpp}       & $\bm{\kappa}(t)$ (Section~\ref{sec:flow}) \\
\code{src/external\_force.cpp} & user-applied bead forces (Section~\ref{sec:external}) \\
\code{src/rng.cpp}        & \code{ran\_1} generator (Chapter~\ref{ch:rng}) \\
\code{python/bdsim/initial.py} & initial configurations (Chapter~\ref{ch:initial}) \\
\code{python/bdsim/properties.py} & observables (Chapter~\ref{ch:observables}) \\
\code{python/bdsim/coarse\_grain.py} & wormlike-chain coarse-graining (Chapter~\ref{ch:coarse}) \\
\code{python/bdsim/dynamics.py} & hydrodynamic radius, $h^{*}$, real and rodlike units (Chapter~\ref{ch:dynamics}) \\
\code{python/bdsim/statistics.py} & correlated-data error bars (Chapter~\ref{ch:statistics}) \\
\bottomrule
\end{longtable}

% ===========================================================================
\chapter{Governing equations}
\label{ch:governing}

\section{The stochastic differential equation}

The configurational distribution function obeys a Fokker--Planck equation whose
It\^{o}-equivalent stochastic differential equation, in the non-dimensional form
integrated by the code, is
\begin{equation}
  \boxed{\;
  \mathrm{d}\R = \left[\bm{K}\cdot\R + \tfrac{1}{4}\,\D\cdot\bm{F}^{\phi}\right]\mathrm{d}t^{*}
                 + \tfrac{1}{\sqrt{2}}\,\B\cdot\mathrm{d}\bm{W}\;}
  \label{eq:sde}
\end{equation}
where $\bm{F}^{\phi}$ is the total conservative force on each bead, $\D$ is the
diffusion tensor, $\B$ satisfies $\D = \B\cdot\B^{\mathsf{T}}$, and
$\mathrm{d}\bm{W}$ is a Wiener increment of variance $\mathrm{d}t^{*}$.

The total force is split into a spring part and everything else,
\begin{equation}
  \bm{F}^{\phi} = \bm{F}^{\mathrm{S}} + \bm{F}^{\mathrm{E}},
  \qquad
  \bm{F}^{\mathrm{E}} = \bm{F}^{\mathrm{EV}} + \bm{F}^{\mathrm{b}} + \bm{F}^{\mathrm{ext}}(t),
\end{equation}
because the scheme of Chapter~\ref{ch:integrator} treats $\bm{F}^{\mathrm{S}}$
implicitly and $\bm{F}^{\mathrm{E}}$ explicitly. In the code $\bm{F}^{\mathrm{E}}$
is assembled by \code{PhysicalModel::non\_spring\_force}, which is the single
place where a new explicit force type would be added.

\section{Flow field}
\label{sec:flow}

The flow enters only through the velocity-gradient tensor $\bm{\kappa}$, applied
to each bead independently:
\begin{equation}
  (\bm{K}\cdot\R)_\nu = \bm{\kappa}(t)\cdot\bm{r}_\nu .
\end{equation}
The code takes $\bm{\kappa}$ directly as a $3\times3$ tensor rather than
enumerating named flow types. It is either constant, or supplied as a table of
samples $(t_i, \bm{\kappa}_i)$ with strictly ascending $t_i$, in which case it is
interpolated componentwise and linearly,
\begin{equation}
  \bm{\kappa}(t) = \bm{\kappa}_i
    + \frac{t-t_i}{t_{i+1}-t_i}\left(\bm{\kappa}_{i+1}-\bm{\kappa}_i\right),
    \qquad t_i \le t < t_{i+1},
\end{equation}
and clamped to the end samples outside the tabulated range. The usual flows are
then just particular tensors, supplied by \code{python/bdsim/flows.py}:
\begin{equation}
  \text{simple shear:}\;\; \kappa_{xy} = \dot\gamma;
  \qquad
  \text{uniaxial extension:}\;\; \bm{\kappa} = \dot\varepsilon\,\mathrm{diag}(1,-\tfrac12,-\tfrac12);
\end{equation}
\begin{equation}
  \text{planar extension:}\;\; \bm{\kappa} = \dot\varepsilon\,\mathrm{diag}(1,-1,0).
\end{equation}
Harmonic traps are not implemented.


\section{External forces}
\label{sec:external}

An external force is not part of the chain's own physics: it is whatever the
experiment applies --- an optical trap, a magnetic bead, a pulling protocol. It is
therefore kept separate from the conservative intramolecular forces and added to
the explicit part of the step, exactly like excluded volume or bending, but
depending on time rather than on configuration.

Each entry applies a force to one bead, either constant or given as a table of
$(t_i, \bm{F}_i)$ samples that is linearly interpolated and clamped outside the
tabulated range --- the same convention as the flow tensor, so a pulling protocol
and a flow protocol are specified in the same way. Negative bead indices count
from the end of the chain, so $-1$ is the last bead.

The canonical use is a stretching experiment, in which equal and opposite forces
are applied to the two ends,
\begin{equation}
  \bm{F}^{\mathrm{ext}}_{1} = -f\hat{\bm{z}},
  \qquad
  \bm{F}^{\mathrm{ext}}_{N} = +f\hat{\bm{z}},
  \label{eq:stretch}
\end{equation}
so that the net force vanishes and the chain is stretched without drifting. This
is what \code{ExternalForce::add\_stretch} builds.

\paragraph{What the external force is deliberately not part of.} It is excluded
from \code{PhysicalModel::non\_spring\_force}, and therefore from the
\code{total\_force} exposed to Python, which continues to mean \emph{the chain's
own force}. Two things depend on that: the Kramers--Kirkwood stress
\eqref{eq:kramers}, which is a sum over intramolecular forces and would be wrong
if an applied load were folded in, and the bit-exact cross-check against the
Fortran (Section~\ref{sec:fortran-compare}), whose force output has no such term.
The applied force is inspected separately through \code{external\_force(phys,
nbeads, t)}.

% ===========================================================================
\chapter{Time integration: the semi-implicit predictor--corrector}
\label{ch:integrator}

The scheme follows Prabhakar and Prakash, with the spring force treated
implicitly so that finitely extensible springs remain stable at moderate
timesteps. It has three stages per step.

\section{Predictor}

An explicit Euler step gives the predicted positions,
\begin{equation}
  \tilde{\R}_{n+1} = \R_n
    + \left[\bm{K}\cdot\R_n
            + \tfrac{1}{4}\D_n\cdot\bm{F}^{\mathrm{S}}_n
            + \tfrac{1}{4}\D_n\cdot\bm{F}^{\mathrm{E}}_n\right]\dt
    + \tfrac{1}{\sqrt{2}}\,\Delta\bm{S}_n ,
  \label{eq:predictor}
\end{equation}
where
\begin{equation}
  \Delta\bm{S}_n = \B_n\cdot\Delta\bm{W}_n
\end{equation}
is the Brownian displacement, constructed once per step
(Section~\ref{sec:brownian-step}) and reused unchanged by the corrector.

\section{Corrector and the \texorpdfstring{$\bm{\Upsilon}$}{Upsilon} term}

The corrector is
\begin{equation}
\begin{aligned}
  \R_{n+1} = \R_n
    &+ \Big[\tfrac{1}{2}\left\{\bm{K}\cdot\R_n + \bm{K}\cdot\tilde{\R}_{n+1}\right\}
       + \tfrac{1}{8}\D_n\cdot\left\{\bm{F}^{\mathrm{S}}_n + \bm{F}^{\mathrm{S}}_{n+1}\right\}\\
    &\phantom{+\Big[}
       + \tfrac{1}{8}\D_n\cdot\left\{\bm{F}^{\mathrm{E}}_n + \tilde{\bm{F}}^{\mathrm{E}}_{n+1}\right\}\Big]\dt
       + \tfrac{1}{\sqrt{2}}\,\Delta\bm{S}_n ,
\end{aligned}
  \label{eq:corrector}
\end{equation}
in which $\tilde{\bm{F}}^{\mathrm{E}}_{n+1}$ is evaluated at the \emph{predicted}
positions, while $\bm{F}^{\mathrm{S}}_{n+1}$ --- the unknown --- is implicit. Any
external force is part of $\bm{F}^{\mathrm{E}}$ and, being a function of time
rather than configuration, is evaluated at $t$ in the predictor and at $t+\dt$ in
the corrector.
Note that $\D_n$ is \emph{not} recomputed at the predicted configuration: the
change over one timestep is negligible and recomputing it was found to make no
difference to accuracy, while costing a second tensor construction.

Collecting everything on the right-hand side of \eqref{eq:corrector} except the
implicit term defines
\begin{equation}
  \tilde{\bm{\Upsilon}}_{n+1} = \R_n
    + \left[\tfrac{1}{2}\left\{\bm{K}\cdot\R_n + \bm{K}\cdot\tilde{\R}_{n+1}\right\}
      + \tfrac{1}{8}\D_n\cdot\bm{F}^{\mathrm{S}}_n
      + \tfrac{1}{8}\D_n\cdot\left\{\bm{F}^{\mathrm{E}}_n+\tilde{\bm{F}}^{\mathrm{E}}_{n+1}\right\}\right]\dt
    + \tfrac{1}{\sqrt{2}}\Delta\bm{S}_n .
  \label{eq:upsilon}
\end{equation}
In the code this is the \code{upsilon} stage, returning a field of $N$ vectors.

\section{Implicit solve for the connector vectors}

Introduce the bond-difference operator
\begin{equation}
  \mathcal{D}_\mu[\bm{X}] \equiv \bm{X}_{\mu+1} - \bm{X}_{\mu}.
\end{equation}
Differencing \eqref{eq:corrector} along the chain and moving the spring term to
the left gives the fixed-point equation solved each step,
\begin{equation}
  \boxed{\;
  \Q^{(j)}_{\mu,n+1} + \frac{\dt}{4}\,\bm{F}^{\mathrm{c}(j)}_{\mu}
    = \bm{\Gamma}^{(j-1)}_{\mu,n+1}\;}
  \label{eq:implicit}
\end{equation}
with
\begin{equation}
  \bm{\Gamma}^{(j-1)}_{\mu,n+1}
    = \mathcal{D}_\mu\!\left[\tilde{\bm{\Upsilon}}_{n+1}\right]
    + \tfrac{1}{8}\,\mathcal{D}_\mu\!\left[\D_n\cdot\bm{F}^{\mathrm{S}(j-1)}_{n+1}\right]\dt
    + \tfrac{1}{4}\,\bm{F}^{\mathrm{c}(j-1)}_{\mu}\dt .
  \label{eq:gamma}
\end{equation}
Here $j$ is the iteration index; $\D_n$ and $\tilde{\bm{\Upsilon}}_{n+1}$ are
fixed during the iteration. Because $\bm{F}^{\mathrm{c}}_\mu$ is parallel to
$\Q_\mu$, equation~\eqref{eq:implicit} implies
\begin{equation}
  \hat{\Q}_{\mu} = \hat{\bm{\Gamma}}_{\mu},
\end{equation}
i.e.\ \emph{the updated bond is parallel to $\bm{\Gamma}$}, so only its magnitude
remains to be found. This is the property that reduces the vector equation to a
scalar one, and it is verified directly by the unit test \code{test\_direction}.

Writing a general force law as
\begin{equation}
  \bm{F}^{\mathrm{c}}(\Q) = H\,\Q\,f(Q),
  \label{eq:general-force-law}
\end{equation}
the magnitude equation is
\begin{equation}
  \left(1 + \frac{\dt}{4}f(Q)\right) Q = \Gamma,
  \qquad \Gamma \equiv \|\bm{\Gamma}_\mu\| ,
  \label{eq:magnitude}
\end{equation}
whose solution for each force law is the subject of Chapter~\ref{ch:springs}.
The bond is then reconstructed as $\Q_\mu = Q\,\bm{\Gamma}_\mu/\Gamma$ and the
bead positions rebuilt by cumulative summation from the (fixed) first bead
$\tilde{\bm{r}}_{1,n+1}$.

\section{Convergence criterion and iteration limit}

The iteration terminates when the root-mean-square change in bead positions
falls below a tolerance,
\begin{equation}
  \frac{1}{N}\left\|\R^{(j)}_{n+1} - \R^{(j-1)}_{n+1}\right\|_2
    < \varepsilon_{\mathrm{tol}},
  \label{eq:convergence}
\end{equation}
where $\|\cdot\|_2$ is the Euclidean norm over all $3N$ components and
$\varepsilon_{\mathrm{tol}}$ is \code{SimParams::implicit\_loop\_tol}. The
iteration is additionally capped at $100N$ iterations. There is no universally
agreed criterion in the literature; \eqref{eq:convergence} is an absolute rather
than relative measure, which is adequate here because the bead coordinates are
$O(1)$ in Hookean units.

\section{Reduced-variance Brownian noise}
\label{sec:noise}

The Wiener increment is not sampled from a true Gaussian. Following \"{O}ttinger,
each component is drawn as
\begin{equation}
  X = \left(Y-\tfrac12\right)\left[c_1\left(Y-\tfrac12\right)^2 + c_2\right],
  \qquad Y \sim \mathcal{U}[0,1],
  \label{eq:ottinger-noise}
\end{equation}
with
\begin{equation}
  c_1 = 14.14855378, \qquad c_2 = 1.21569221 .
\end{equation}
This reproduces the correct second and fourth moments, which is sufficient for
averages obtained from Brownian dynamics, at appreciably lower cost than a
Gaussian generator. The increment used in \eqref{eq:predictor} is then
\begin{equation}
  \Delta\bm{S}_n = \sqrt{\dt}\;\B_n\cdot\bm{X}_0 ,
\end{equation}
with $\bm{X}_0$ the field of deviates \eqref{eq:ottinger-noise}.

\section{Stepping, sampling and reproducibility}

The integrator advances a chain over a requested time span and nothing else:
sampling, property evaluation and output are owned by the Python driver, which
integrates in segments. The random stream persists across calls, so a segmented
run is identical to one continuous run --- this is what makes it legitimate to
write a snapshot every $N_{\mathrm{write}}$ steps without perturbing the
trajectory. Trajectories in an ensemble are seeded by index, so ensemble results
do not depend on how the work is distributed over workers.

% ===========================================================================
\chapter{Spring force laws and the implicit solve}
\label{ch:springs}

\section{Force laws}

Every law is expressed in the form \eqref{eq:general-force-law},
$\bm{F}^{\mathrm{c}} = H\Q f(Q)$. Internally the code works with the
\emph{reduced} bond length and reduced natural length
\begin{equation}
  q = \frac{Q}{\sqrt{b}}, \qquad q_0 = \frac{\sigma}{\sqrt{b}},
\end{equation}
where $\sqrt{b}$ is the finite-extensibility parameter (\code{SpringParams::sqrtb})
and $\sigma$ the natural length (\code{SpringParams::natural\_length}). The
implemented $f$ are listed in Table~\ref{tab:force-laws}.

\begin{table}[ht]
\centering
\caption{Spring force laws, as implemented in \code{force\_sans\_hookean}.
$q=Q/\sqrt{b}$ is the reduced bond length and $q_0=\sigma/\sqrt{b}$ the reduced
natural length.}
\label{tab:force-laws}
\begin{tabular}{@{}lll@{}}
\toprule
Law & \code{Spring} & $f(q)$ \\
\midrule
Hookean & \code{Hook} & $1$ \\[0.4em]
FENE & \code{FENE} & $\dfrac{1}{1-q^2}$ \\[0.7em]
Inverse Langevin (Pad\'e) & \code{ILC} & $\dfrac{3-q^2}{3\left(1-q^2\right)}$ \\[0.7em]
Worm-like chain (Marko--Siggia) & \code{WLC} &
  $\dfrac{1}{6q}\left[4q + \dfrac{1}{(1-q)^2} - 1\right]$ \\[0.7em]
Fraenkel & \code{Fraenkel} & $1 - \dfrac{q_0}{q}$ \\[0.7em]
FENE--Fraenkel & \code{FENEFraenkel} &
  $\dfrac{1-q_0/q}{1-(q-q_0)^2}$ \\[0.7em]
Bounded WLC & \code{WLCbounded} & see \eqref{eq:wlcb-f} \\
\bottomrule
\end{tabular}
\end{table}

The bounded (modified Marko--Siggia) law uses the reduced variable
$t = (q-q_0)/(1-q_0)$ and reads
\begin{equation}
  f(q) = \frac{2}{3q}\left[
    \frac{(1-t)^{-2}-1}{4} + t
    - q_0\frac{(1+t)^{-2}-1}{4} + q_0 t
  \right].
  \label{eq:wlcb-f}
\end{equation}

Two limits are worth recording because they are asserted by the unit tests:
FENE--Fraenkel with $q_0 = 0$ reduces exactly to FENE (so $\alpha \equiv b$ plays
the role of the FENE $b$ parameter), and bounded WLC with $q_0 = 0$ reduces
exactly to Marko--Siggia WLC.

\section{Bead forces from bond forces}

The force carried by bond $\mu$ is $\bm{F}^{\mathrm{c}}_\mu = f(q_\mu)\,\Q_\mu$
(with $H=1$), and the force on the beads follows by differencing,
\begin{equation}
  \bm{F}^{\mathrm{S}}_1 = \bm{F}^{\mathrm{c}}_1, \qquad
  \bm{F}^{\mathrm{S}}_\nu = \bm{F}^{\mathrm{c}}_\nu - \bm{F}^{\mathrm{c}}_{\nu-1}
    \;\; (1<\nu<N), \qquad
  \bm{F}^{\mathrm{S}}_N = -\bm{F}^{\mathrm{c}}_{N-1}.
\end{equation}

\section{Solving the implicit equation}

Equation~\eqref{eq:magnitude} must be solved for $q$ given
$\gamma \equiv \Gamma/\sqrt{b}$ and $\dt/4$. The laws fall into three families,
handled by three different techniques; \code{solve\_implicit\_r} dispatches on
the family and returns both the root and $f$ evaluated at it.

\subsection{Closed form}

\emph{Hookean:}
\begin{equation}
  q = \frac{\gamma}{1 + \dt/4}, \qquad f = 1 .
\end{equation}

\emph{Fraenkel:}
\begin{equation}
  q = \frac{\gamma + (\dt/4)\,q_0}{1 + \dt/4}, \qquad f = 1 - \frac{q_0}{q}.
\end{equation}

\subsection{Cubic family}

For FENE, ILC, WLC and FENE--Fraenkel, \eqref{eq:magnitude} rearranges into a
cubic
\begin{equation}
  c_3 q^3 + c_2 q^2 + c_1 q + c_0 = 0, \qquad c_3 = 1,
\end{equation}
with the coefficients given in Table~\ref{tab:cubic-coeffs}
(writing $\beta \equiv \dt/4$).

\begin{table}[ht]
\centering
\caption{Cubic coefficients, with $\beta = \dt/4$, $\gamma$ the reduced
right-hand side and $q_0$ the reduced natural length
(\code{cubic\_coefficients}).}
\label{tab:cubic-coeffs}
\renewcommand{\arraystretch}{1.5}
\begin{tabular}{@{}llll@{}}
\toprule
Law & $c_0$ & $c_1$ & $c_2$ \\
\midrule
FENE & $\gamma$ & $-(1+\beta)$ & $-\gamma$ \\
ILC  & $\dfrac{3\gamma}{3+\beta}$ & $-\dfrac{3+3\beta}{3+\beta}$ & $-\dfrac{3\gamma}{3+\beta}$ \\
WLC  & $-\dfrac{3\gamma}{d}$ & $\dfrac{3(1+\beta+2\gamma)}{d}$ & $-\dfrac{3(4+3\beta+2\gamma)}{2d}$ \\
FENE--Fraenkel & $\beta q_0 + \gamma(1-q_0^2)$ & $-1+q_0^2-\beta+2\gamma q_0$ & $-(2q_0+\gamma)$ \\
\bottomrule
\end{tabular}

\vspace{0.5em}
{\footnotesize For WLC, $d = 2\left(\tfrac{3}{2}+\beta\right)$.}
\end{table}

\subsubsection{Root selection for FENE--Fraenkel}

The FENE--Fraenkel cubic is solved by the direct (trigonometric/Cardano) formula,
selecting the root that lies in the physical interval. That such a root exists,
and is unique, follows from evaluating the cubic $g(q)$ at the boundaries. In
dimensional form with $\alpha = \delta Q^2$ and $\beta' = \tfrac14\alpha\dt$,
\begin{align}
  g(\infty) &> 0, \\
  g(\sigma+\sqrt{\alpha}) &= -\beta'\sqrt{\alpha} < 0, \\
  g(\sigma-\sqrt{\alpha}) &= +\beta'\sqrt{\alpha} > 0, \\
  g(0) &= \beta' + (\alpha-\sigma^2)\,\Gamma > 0 \quad\text{if } \sqrt{\alpha}\ge\sigma, \\
  g(-\infty) &< 0 .
\end{align}
Hence exactly one root lies between $\max(0,\,\sigma-\sqrt{\alpha})$ and
$\sigma+\sqrt{\alpha}$, and that is the physical bond length. In reduced units
the bracket used by the code is
\begin{equation}
  q \in \left[\max(q_0-1,\,0),\; q_0+1\right].
  \label{eq:ff-bracket}
\end{equation}

\subsubsection{Approximate root plus Newton polish}

For FENE, ILC and WLC the root is instead approximated and then polished by
Newton--Raphson on the cubic. The initial guess is Hookean at small $\gamma$ and
asymptotic at large $\gamma$ (where $q\to1$):
\begin{equation}
  q^{(0)} =
  \begin{cases}
    \dfrac{\gamma}{1+\beta}, & \gamma < 1,\\[0.8em]
    1 - \dfrac{\beta}{2\gamma}, & \text{FENE},\\[0.7em]
    1 - \dfrac{\beta}{3\gamma}, & \text{ILC},\\[0.7em]
    1 - \sqrt{\dfrac{\beta}{6\gamma}}, & \text{WLC}.
  \end{cases}
\end{equation}
The polish is skipped for WLC when $\gamma > 100$, where Newton iterations
wander off the physical root and the asymptotic guess is already accurate.

\subsection{Transcendental family: bounded WLC}

The bounded WLC does not reduce to a polynomial. Its implicit equation is
\begin{multline}
  \frac{6\gamma}{\dt} = \frac{6q}{\dt}
    + \frac{1}{4}\left(\frac{1-\sigma}{1-q}\right)^{2} - \frac{1}{4}
    + \frac{q-\sigma}{1-\sigma} \\
    - \frac{\sigma}{4}\left(\frac{1-\sigma}{1+q-2\sigma}\right)^{2}
    + \frac{\sigma}{4} + \sigma\frac{q-\sigma}{1-\sigma},
  \label{eq:wlcb-implicit}
\end{multline}
which is solved by \code{rtsafe} (safeguarded Newton with bisection fallback) on
the bracket
\begin{equation}
  q \in \left[\max(0,\,2\sigma-1)+\epsilon,\; 1-\epsilon\right],
\end{equation}
$\epsilon$ being machine epsilon. The analytic derivative required by
\code{rtsafe} is
\begin{equation}
  \frac{\partial}{\partial q} = \frac{6}{\dt} + \frac{1}{1-\sigma} + \frac{\sigma}{1-\sigma}
    + \frac{(1-\sigma)^2}{2(1-q)^3}
    + \frac{(1-\sigma)^2\sigma}{2(1-2\sigma+q)^3}.
\end{equation}

\paragraph{Lookup tables are not implemented.} The thesis describes an adaptive
lookup table built once at start-up to avoid repeated \code{rtsafe} calls for
this law. The current code solves \eqref{eq:wlcb-implicit} directly at every
step. This is a deliberate simplification: it removes a large amount of
table-construction machinery and a class of interpolation-accuracy questions, at
the cost of making the bounded WLC the slowest law. If bounded-WLC runs become a
bottleneck the table can be reinstated behind the same interface, since
\code{solve\_implicit\_r} is a pure function of $(\text{law}, \beta, \gamma, q_0)$.

\section{Robustness of the solve}
\label{sec:spring-robustness}

Two safeguards are present that are not in the thesis. Both were added after
diagnosing intermittent blow-ups in long FENE--Fraenkel runs with HI, and both
are covered by \code{tests/test\_hi\_stability.cpp}.

\paragraph{Bracketed fallback for the cubic.} The closed-form cubic solution
loses accuracy for extreme coefficients. For $\gamma \gtrsim 10^{10}$ --- which
can occur transiently if the Brownian displacement is momentarily very large ---
the trigonometric formula could return a spurious root of order $\gamma$,
\emph{outside} the physical bracket \eqref{eq:ff-bracket}, bypassing finite
extensibility entirely and producing an unbounded bond. \code{find\_roots\_cubic}
now checks that the returned root lies in $[\,q_{\mathrm{lo}}, q_{\mathrm{hi}}]$
and, if not, falls back to bisection on that bracket, which cannot fail.

\paragraph{Strict interior clamping.} After polishing, the root is clamped
strictly inside the singularities of $f$:
\begin{equation}
  q \leftarrow \min\!\left(\max(q,\; q_{\mathrm{lo}}+\delta),\; q_{\mathrm{hi}}-\delta\right),
  \qquad \delta = 10^{-6},
\end{equation}
with $(q_{\mathrm{lo}},q_{\mathrm{hi}}) = (\max(q_0-1,0),\,q_0+1)$ for
FENE--Fraenkel and $(0,1)$ for FENE/ILC/WLC. Because $f$ diverges at those
endpoints, landing exactly on one produces a non-finite force; the clamp
guarantees the spring force is finite for \emph{any} $\gamma$, which is precisely
the stability property the semi-implicit scheme is supposed to provide. Tests
assert boundedness up to $\gamma = 10^{20}$.

% ===========================================================================
\chapter{Hydrodynamic interaction}
\label{ch:hi}

Setting $h^{*} = 0$ gives free draining, $\D = \bm{\delta}$ (identity), and
$\B = \bm{\delta}$; the Brownian displacement is then simply
$\sqrt{\dt}\,\bm{X}_0$. Everything below concerns $h^{*} > 0$.

\section{The RPY diffusion tensor}
\label{sec:rpy}

$\D$ is a $3N\times3N$ block matrix. The diagonal blocks are the identity (bead
self-mobility), and the off-diagonal blocks use the Rotne--Prager--Yamakawa form,
which is constructed to remain positive semi-definite for all configurations.
For beads $\mu\ne\nu$ separated by $\bm{b} = \bm{r}_\nu - \bm{r}_\mu$ with
$r = \|\bm{b}\|$,
\begin{equation}
  \D_{\mu\nu} = \omega_1(r)\,\bm{\delta} + \omega_2(r)\,\bm{b}\bm{b},
  \label{eq:rpy-block}
\end{equation}
with two branches at the overlap radius $r = 2a^{*} = 2\sqrt{\pi}h^{*}$:
\begin{equation}
  \boxed{
  \begin{aligned}
    r \ge 2\sqrt{\pi}h^{*}:\quad
      \omega_1 &= \tfrac{3}{4}\sqrt{\pi}\,\frac{h^{*}}{r}
                  \left[1 + \tfrac{2\pi}{3}\left(\tfrac{h^{*}}{r}\right)^{2}\right], &
      \omega_2 &= \tfrac{3}{4}\sqrt{\pi}\,\frac{h^{*}}{r^{3}}
                  \left[1 - 2\pi\left(\tfrac{h^{*}}{r}\right)^{2}\right],\\[0.5em]
    r < 2\sqrt{\pi}h^{*}:\quad
      \omega_1 &= 1 - \frac{9}{32\sqrt{\pi}}\,\frac{r}{h^{*}}, &
      \omega_2 &= \frac{3}{32\sqrt{\pi}}\,\frac{1}{h^{*} r} .
  \end{aligned}}
  \label{eq:rpy-omegas}
\end{equation}
Substituting $a^{*} = \sqrt{\pi}h^{*}$ recovers the familiar textbook form: the
far-field branch is
$\D_{\mu\nu} = \tfrac{3a}{4r}\left[\left(1+\tfrac{2a^{2}}{3r^{2}}\right)\bm{\delta}
 + \left(1-\tfrac{2a^{2}}{r^{2}}\right)\tfrac{\bm{b}\bm{b}}{r^{2}}\right]$
and the near-field branch is
$\D_{\mu\nu} = \left(1-\tfrac{9r}{32a}\right)\bm{\delta} + \tfrac{3}{32ar}\bm{b}\bm{b}$.
Squared separations below $10^{-12}$ are clamped to avoid a singularity for
coincident beads.

\section{The Brownian displacement}
\label{sec:brownian-step}

What is actually required is never $\B$ itself but the product
\begin{equation}
  \Delta\bm{S} = \sqrt{\dt}\;\B\cdot\bm{X}_0,
  \qquad \D = \B\cdot\B^{\mathsf{T}} .
  \label{eq:delS}
\end{equation}
Note that \eqref{eq:delS} does not determine $\B$ uniquely --- any $\B$
satisfying $\D = \B\B^{\mathsf{T}}$ gives $\Delta\bm{S}$ the correct covariance,
which is all the SDE requires. The code offers two constructions.

\subsection{Cholesky factorisation}

$\B = \bm{L}$, the lower triangular Cholesky factor of $\D$, so that
$\Delta S_i = \sum_{j\le i} L_{ij} X_{0,j}$. This is exact and unconditionally
stable, and costs $O(N^3)$. It is the reference against which the approximate
method is checked.

\subsection{Chebyshev approximation}

The alternative approximates the \emph{symmetric} square root by a Chebyshev
polynomial in $\D$, requiring only matrix--vector products and hence
$O(N_T N^2)$ work for $N_T$ terms.

The construction rests on two facts: the eigenvectors of $\D$ are also
eigenvectors of any polynomial in $\D$, with eigenvalues transformed by that
polynomial; and $\sqrt{\D} = \sum_k \lambda_k^{1/2}\bm{v}_k\bm{v}_k$ in the
spectral representation. So a scalar polynomial approximating $\sqrt{\lambda}$
over the spectral range, applied to $\D$, approximates $\sqrt{\D}$.

Let the spectrum lie in $[\lambda_{\min}, \lambda_{\max}]$. Map $\D$ onto the
Chebyshev interval $[-1,1]$ via
\begin{equation}
  \D' = d_a \D + d_b \bm{\delta}, \qquad
  d_a = \frac{2}{\lambda_{\max}-\lambda_{\min}}, \qquad
  d_b = -\frac{\lambda_{\max}+\lambda_{\min}}{\lambda_{\max}-\lambda_{\min}} .
\end{equation}
The coefficients for $\sqrt{\cdot}$ with $L+1$ nodes are
\begin{equation}
  \lambda_k = \frac{1}{d_a}\cos\!\left[\frac{\pi(k+\tfrac12)}{L+1}\right] - \frac{d_b}{d_a},
  \qquad k = 0,\dots,L,
\end{equation}
\begin{equation}
  a_p = \frac{2}{L+1}\sum_{k=0}^{L}\sqrt{\lambda_k}\,
        \cos\!\left[\frac{p\,(k+\tfrac12)\pi}{L+1}\right],
  \qquad a_0 \leftarrow \tfrac{a_0}{2}.
\end{equation}
The product with $\bm{X}_0$ is then accumulated using the three-term recurrence
directly on vectors, never forming a matrix:
\begin{equation}
  \bm{V}_0 = \bm{X}_0, \qquad
  \bm{V}_1 = \D'\cdot\bm{X}_0, \qquad
  \bm{V}_p = 2\,\D'\cdot\bm{V}_{p-1} - \bm{V}_{p-2},
\end{equation}
\begin{equation}
  \B\cdot\bm{X}_0 \;\approx\; \sum_{p=0}^{L} a_p \bm{V}_p .
\end{equation}
The initial number of terms follows Fixman's scaling,
\begin{equation}
  L = \left\lfloor \mu_{\mathrm{cheb}}\sqrt{\frac{\lambda_{\max}}{\lambda_{\min}}} + \tfrac12 \right\rfloor + 1,
  \label{eq:ncheb}
\end{equation}
clamped to $[2, N_{\mathrm{cheb}}^{\max}]$ with $N_{\mathrm{cheb}}^{\max} = 500$
and $\mu_{\mathrm{cheb}}$ a user multiplier. Since
$\lambda_{\max}/\lambda_{\min}$ scales roughly as $N^{1/2}$ near equilibrium and
as $N$ for stretched chains, the total cost scales between $N^{2.25}$ and
$N^{2.5}$, versus $N^{3}$ for Cholesky.

\subsection{Fluctuation--dissipation error control}

After evaluating the series, the fluctuation--dissipation identity is checked:
\begin{equation}
  \varepsilon_{\mathrm{FD}}
    = \frac{\left|\;\|\B\cdot\bm{X}_0\|^{2} - \bm{X}_0\cdot\D\cdot\bm{X}_0\;\right|}
           {\bm{X}_0\cdot\D\cdot\bm{X}_0} ,
  \label{eq:fd}
\end{equation}
which must vanish for an exact square root. If $\varepsilon_{\mathrm{FD}}$
exceeds the tolerance \code{fd\_err\_max}, $L$ is increased by $\max(1, L/4)$ and
the series re-evaluated, up to $N_{\mathrm{cheb}}^{\max}$.

\section{Spectral bounds}
\label{sec:spectral-bounds}

\subsection{Why the interval must contain the spectrum}

The Chebyshev polynomial approximates $\sqrt{\cdot}$ \emph{only} on
$[\lambda_{\min},\lambda_{\max}]$. Outside that interval it does not merely lose
accuracy: Chebyshev polynomials grow like $\cosh$ beyond $\pm1$, so any
eigencomponent of $\bm{X}_0$ whose eigenvalue lies below $\lambda_{\min}$ (or
above $\lambda_{\max}$) is \emph{amplified} rather than square-rooted. Repeated
over many steps this pumps energy into the corresponding mode and the trajectory
diverges. Containment of the spectrum is therefore a correctness requirement, not
an accuracy preference.

\subsection{The Fixman estimate and why it was replaced}

The thesis (and the original Fortran) estimates the extremes with two
matrix--vector products,
\begin{equation}
  d_{\max} = \frac{2}{3N}\,\D : \bm{u}^{+}\bm{u}^{+},
  \qquad
  d_{\min} = \frac{1}{2\cdot 3N}\,\D : \bm{u}^{-}\bm{u}^{-},
  \label{eq:fixman}
\end{equation}
where $\bm{u}^{+}$ has all components $+1$ (fully cooperative motion, fastest
diffusion) and $\bm{u}^{-}$ alternates in sign (the opposite extreme). This is
cheap and physically motivated.

It is, however, \emph{not a bound}. Both expressions are Rayleigh quotients (up
to constant factors), and a Rayleigh quotient always lies strictly between
$\lambda_{\min}$ and $\lambda_{\max}$. Consequently \eqref{eq:fixman}
systematically \emph{overestimates} $\lambda_{\min}$. For compact configurations
the gap is harmless, but stretched chains develop a small eigenvalue that
$\bm{u}^{-}$ does not resolve. Measured on a stretched $N=30$ configuration, the
estimate gave $d_{\min} = 0.379$ against a true $\lambda_{\min} = 0.139$; noise
aligned with the smallest eigenvector then produced
$\varepsilon_{\mathrm{FD}} \sim 10^{158}$ and an immediate blow-up. This was the
root cause of the intermittent instabilities discussed in
Section~\ref{sec:spring-robustness}.

\subsection{Lanczos estimate with Sturm bisection (current implementation)}

The code now estimates the extreme eigenvalues by an $m$-step Lanczos iteration
with full re-orthogonalisation ($m = 30$ by default, capped at $3N$). Lanczos
builds an orthonormal basis $\{\bm{v}_1,\dots,\bm{v}_m\}$ and a symmetric
tridiagonal projection
\begin{equation}
  \bm{T}_m =
  \begin{pmatrix}
    \alpha_1 & \beta_1 & & \\
    \beta_1 & \alpha_2 & \ddots & \\
     & \ddots & \ddots & \beta_{m-1}\\
     & & \beta_{m-1} & \alpha_m
  \end{pmatrix},
  \qquad
  \begin{aligned}
    \alpha_j &= \bm{v}_j\cdot\D\cdot\bm{v}_j,\\
    \beta_j  &= \|\bm{w}_j\|,
  \end{aligned}
\end{equation}
whose extreme eigenvalues (Ritz values) converge to those of $\D$ far faster than
power iteration. The extreme Ritz values are extracted by bisection on the Sturm
sequence of $\bm{T}_m$: the number of eigenvalues of $\bm{T}_m$ below $x$ is the
number of negative entries in
\begin{equation}
  d_1 = \alpha_1 - x, \qquad
  d_j = (\alpha_j - x) - \frac{\beta_{j-1}^{2}}{d_{j-1}},
\end{equation}
and the search interval is initialised from the Gershgorin bounds of
$\bm{T}_m$.

Ritz values bracket the true spectrum from the \emph{inside}, so they are padded
outwards to obtain the interval actually used:
\begin{equation}
  \lambda_{\min} = 0.5\,\theta_{\min}, \qquad
  \lambda_{\max} = 1.02\,\theta_{\max},
  \label{eq:padding}
\end{equation}
with $\theta$ the extreme Ritz values. Verified over $200$ configurations sampled
along real trajectories, \eqref{eq:padding} contained the spectrum in every case.

\subsection{Warm-starting the Lanczos iteration}
\label{sec:warm-start}

Lanczos converges to the extremes quickly, but $m$ matrix--vector products per
step is still expensive next to the two of \eqref{eq:fixman}. The saving comes
from the observation that over a single timestep the configuration --- and hence
the spectrum \emph{and its extremal eigenvectors} --- barely changes. Seeding the
Krylov space with the previous step's extremal Ritz vectors therefore reaches the
same accuracy in far fewer iterations: the implementation uses $m = 30$ for a cold
start and $m = 8$ once warm.

The Ritz vectors are recovered at no extra matrix--vector cost. The Lanczos basis
$\bm{V}$ is already stored for re-orthogonalisation, and the eigenvector $\bm{y}$
of the tridiagonal $\bm{T}_m$ for a converged Ritz value is obtained by two steps
of inverse iteration (an $O(m)$ tridiagonal solve, shifted slightly off the
eigenvalue to stay non-singular); the Ritz vector is then
$\bm{V}^{\mathsf{T}}\bm{y}$.

Two failure modes are guarded explicitly.
\begin{itemize}
  \item \emph{Krylov breakdown.} Seeding with $\bm{v}_{\min}$ alone is
        dangerous: if it is already close to an exact eigenvector, Lanczos
        terminates after one step and never sees $\lambda_{\max}$. The seed is
        therefore built to be rich in \emph{both} extremes,
        \begin{equation}
          \bm{v}_{\mathrm{seed}} \propto
            \bm{v}_{\min} + \bm{v}_{\max} + 10^{-2}\,\hat{\bm{u}},
        \end{equation}
        with $\hat{\bm{u}}$ the deterministic cold-start vector, whose small
        admixture makes degeneracy impossible.
  \item \emph{Early termination.} If a warm run still collapses into an invariant
        subspace in fewer than four iterations, its bounds are discarded and the
        step is redone cold.
\end{itemize}

Correctness does not rest on any of this. The bounds feed only the choice of
Chebyshev interval; the fluctuation--dissipation test \eqref{eq:fd} and the
Cholesky fallback below still police the result, so a poor seed costs time, never
accuracy. That the warm-started interval nevertheless does contain the spectrum is
asserted directly by \code{test\_hi\_stability}, which compares the bounds against
exact (Jacobi) eigenvalues at intervals along an evolving trajectory.

\subsection{Cholesky fallback}

Padding makes containment overwhelmingly likely but cannot make it a theorem. The
final safety net is therefore in \code{Diffusion::brownian\_step}: if
$\varepsilon_{\mathrm{FD}}$ from \eqref{eq:fd} still exceeds tolerance after the
term-growth loop, that single step falls back to the exact Cholesky factor. This
trades speed for correctness on the rare problematic step and makes divergence
structurally impossible.

\subsection{Cost}

The bound estimator is the dominant cost of the Chebyshev path, so warm-starting
it (Section~\ref{sec:warm-start}) roughly halves the total. Before warm-starting,
at $N=51$, $h^{*}=0.2$, the Chebyshev step cost $433\,\mu$s of which about
$330\,\mu$s ($\approx76\%$) was the estimator; it now costs $223\,\mu$s.
Table~\ref{tab:hi-cost} gives the effect across chain lengths.

The practical consequence is that the crossover between the two square-root
methods has moved down from $N \approx 100$ to $N \approx 40$: at $N=51$
Chebyshev is now the faster method, where previously the exact Cholesky
factorisation beat it.

\begin{table}[ht]
\centering
\caption{Measured per-step cost of the two square-root methods, FENE--Fraenkel
chain, $h^{*}=0.2$, $\dt = 10^{-3}$, single core. ``Chebyshev before'' is with a
cold Lanczos start at every step.}
\label{tab:hi-cost}
\begin{tabular}{@{}rrrrrl@{}}
\toprule
$N$ & Cheb.\ before ($\mu$s) & Cheb.\ now ($\mu$s) & speed-up & Cholesky ($\mu$s) & faster now \\
\midrule
20  & 112  & 37   & $3.1\times$ & 24    & Cholesky \\
51  & 479  & 223  & $2.2\times$ & 250   & Chebyshev \\
100 & 1760 & 1037 & $1.7\times$ & 1983  & Chebyshev \\
200 & 8028 & 4659 & $1.7\times$ & 15799 & Chebyshev \\
\bottomrule
\end{tabular}
\end{table}

\paragraph{The cache is per integration call.} \code{SpectralCache} lives in the
\code{Mobility}, which is constructed by \code{PhysicalModel} once per call to
\code{time\_integrate\_chain}. The warm start therefore accelerates steps
\emph{within} one call but is cold on the first step of each call. This is
immaterial for normal use, where a call spans many steps, but it does mean the
benefit is lost if the integrator is driven one step at a time:

\begin{center}
\begin{tabular}{@{}rr@{}}
\toprule
steps per call & cost ($\mu$s/step, $N=51$) \\
\midrule
200 & 217 \\
50  & 225 \\
10  & 239 \\
1   & 439 \\
\bottomrule
\end{tabular}
\end{center}

Driving the integrator a single step at a time --- as an external ODE coupling
would --- returns the cost to the cold-start value. Recovering it there would mean
giving the Python layer a persistent integrator object that owns the
\code{PhysicalModel} across calls, rather than rebuilding it per call.

\paragraph{Further optimisation.} One further saving remains available and is not
implemented: recomputing the bounds only every $K$ steps and reusing them in
between, which is safe for the same reason warm-starting is (the spectrum moves
slowly) and is likewise policed by the fluctuation--dissipation check.

% ===========================================================================
\chapter{Bending potential}
\label{ch:bending}

\section{Potential and forces}

The implemented bending potential is
\begin{equation}
  \frac{\phi_{\mathrm{b},\mu}}{\kB T} = C\left(1-\cos\theta_\mu\right),
  \label{eq:bending-potential}
\end{equation}
where $C$ is the rigidity constant (\code{BendingParams::stiffness}) and
\begin{equation}
  \theta_\mu = \arccos\left(\uv_{\mu-1}\cdot\uv_{\mu}\right)
\end{equation}
is the angle at bead $\mu$ between adjacent bond unit vectors.

Bead $\mu$ feels contributions from three angles --- its own and its two
neighbours':
\begin{equation}
  \bm{F}^{\mathrm{b}}_\mu
    = \bm{F}^{\mathrm{b}}_{\mu,\mu}
    + \bm{F}^{\mathrm{b}}_{\mu,\mu-1}
    + \bm{F}^{\mathrm{b}}_{\mu,\mu+1},
\end{equation}
each obtained as $-\partial\phi_{\mathrm{b},\mu'}/\partial\bm{r}_\mu$ via the
chain rule. The geometric factors are
\begin{equation}
  \frac{\partial\theta_\mu}{\partial\uv_\mu} = -\frac{\uv_{\mu-1}}{\sin\theta_\mu},
  \qquad
  \frac{\partial\theta_\mu}{\partial\uv_{\mu-1}} = -\frac{\uv_{\mu}}{\sin\theta_\mu},
\end{equation}
\begin{equation}
  \frac{\partial\uv_\mu}{\partial\bm{r}_\mu} = \frac{1}{Q_\mu}\left(\uv_\mu\uv_\mu-\bm{\delta}\right),
  \qquad
  \frac{\partial\uv_{\mu-1}}{\partial\bm{r}_\mu} = \frac{1}{Q_{\mu-1}}\left(\bm{\delta}-\uv_{\mu-1}\uv_{\mu-1}\right).
\end{equation}
Combining these gives the general result
\begin{multline}
  \bm{F}^{\mathrm{b}}_\mu =
   \frac{\partial\phi_{\mathrm{b},\mu}}{\partial\theta_\mu}\frac{1}{\sin\theta_\mu}
   \left[\frac{\cos\theta_\mu\,\uv_\mu-\uv_{\mu-1}}{Q_\mu}
       + \frac{-\cos\theta_\mu\,\uv_{\mu-1}+\uv_\mu}{Q_{\mu-1}}\right]\\
 + \frac{\partial\phi_{\mathrm{b},\mu-1}}{\partial\theta_{\mu-1}}\frac{1}{\sin\theta_{\mu-1}}
   \left[\frac{-\cos\theta_{\mu-1}\,\uv_{\mu-1}+\uv_{\mu-2}}{Q_{\mu-1}}\right]\\
 + \frac{\partial\phi_{\mathrm{b},\mu+1}}{\partial\theta_{\mu+1}}\frac{1}{\sin\theta_{\mu+1}}
   \left[\frac{\cos\theta_{\mu+1}\,\uv_\mu-\uv_{\mu+1}}{Q_\mu}\right].
\end{multline}

For the specific potential \eqref{eq:bending-potential},
$-\partial\phi_{\mathrm{b},i}/\partial\theta_i = -\kB T\,C\sin\theta_i$, so every
$\sin\theta$ cancels --- which is why no small-angle guard is needed in the code
--- leaving
\begin{equation}
  \boxed{
  \begin{aligned}
  \frac{\bm{F}^{\mathrm{b}}_\mu}{\kB T\,C} =\;
   &\left[\frac{\cos\theta_\mu\,\uv_\mu-\uv_{\mu-1}}{Q_\mu}
       + \frac{-\cos\theta_\mu\,\uv_{\mu-1}+\uv_\mu}{Q_{\mu-1}}\right] \\
 &+ \left[\frac{-\cos\theta_{\mu-1}\,\uv_{\mu-1}+\uv_{\mu-2}}{Q_{\mu-1}}\right]
  + \left[\frac{\cos\theta_{\mu+1}\,\uv_\mu-\uv_{\mu+1}}{Q_\mu}\right].
  \end{aligned}}
  \label{eq:bending-force}
\end{equation}
The chain ends keep only the terms that exist: bead $1$ retains only the third
bracket and bead $N$ only the second, which is exactly how \code{bending.cpp}
special-cases $\nu = 1$, $2$, $N-1$, $N$ (and the whole $N=3$ case) around the
general interior loop.

\paragraph{A sign discrepancy in the literature.} Equation~\eqref{eq:bending-force}
differs from the corresponding expression of Saadat and Khomami in the sign of
the $-\cos\theta_{\mu-1}\uv_{\mu-1}$ term. Given that the form above reproduces
hand-computed forces for 3-, 4- and 5-bead configurations and yields the correct
equilibrium angle distribution, the published expression is very likely a
typographical error.

\section{Equilibrium angle distribution}

Because the spring and bending contributions to the configurational integral
separate (the spring potential depends only on $Q$, not orientation), the
equilibrium angle distribution follows from the bending potential alone:
\begin{equation}
  \psi_{\mathrm{eq}}(\theta)
    = \frac{\sin\theta\,e^{-\phi_\mathrm{b}/\kB T}}
           {\int_0^\pi \sin\theta\,e^{-\phi_\mathrm{b}/\kB T}\,\mathrm{d}\theta}
    = \left[\frac{C}{1-e^{-2C}}\right]\sin\theta\;e^{-C(1-\cos\theta)} ,
  \label{eq:angle-dist}
\end{equation}
with moments
\begin{equation}
  \langle\cos\theta\rangle = \coth C - \frac{1}{C},
  \qquad
  \langle\cos^{2}\theta\rangle
    = \frac{2e^{-C}\left[(C^{2}+2)\sinh C - 2C\cosh C\right]}{C^{2}\left(1-e^{-2C}\right)} .
\end{equation}
Equation~\eqref{eq:angle-dist} is both a validation target and the sampling
distribution used to build correlated initial configurations
(Section~\ref{sec:mc-chain}).

% ===========================================================================
\chapter{Excluded volume}
\label{ch:ev}

Excluded volume acts pairwise between beads separated by at least
\code{contour\_dist\_for\_EV} along the chain. For each such pair, with
$\bm{q} = \bm{r}_\nu-\bm{r}_\mu$ and $r = \|\bm{q}\|$, a scalar factor
$\Phi(r)$ is formed and the force applied antisymmetrically,
\begin{equation}
  \bm{F}^{\mathrm{EV}}_\mu \mathrel{-}= \Phi(r)\,\bm{q},
  \qquad
  \bm{F}^{\mathrm{EV}}_\nu \mathrel{+}= \Phi(r)\,\bm{q},
\end{equation}
so that $\Phi > 0$ is repulsive. Three functional forms are implemented.

\paragraph{Gaussian.}
\begin{equation}
  \Phi_{\mathrm{G}}(r) = \frac{z^{*}}{d^{*5}}\exp\!\left(-\frac{r^{2}}{2d^{*2}}\right),
\end{equation}
with $z^{*}$ the dimensionless EV strength and $d^{*}$ the EV range.

\paragraph{Lennard-Jones.}
\begin{equation}
  \Phi_{\mathrm{LJ}}(r) = \frac{48 z^{*} d^{*12}}{r^{14}} - \frac{24 z^{*} d^{*6}}{r^{8}},
\end{equation}
i.e.\ the standard $12$--$6$ force divided by $r$. It is evaluated for
$R_{\min} \le r \le R_{\mathrm{cut}}$, held at its $R_{\min}$ value for
$r < R_{\min}$ (a soft core that prevents overflow at small separations), and set
to zero beyond the cutoff. The cutoff depends on the run mode:
\begin{equation}
  R_{\mathrm{cut}} =
  \begin{cases}
    d^{*}2^{1/6}, & \text{equilibration (purely repulsive, WCA)},\\
    R_{\mathrm{cut}}^{\mathrm{p}}, & \text{production.}
  \end{cases}
\end{equation}

\paragraph{Soddemann--D\"{u}nweg--Kremer (SDK).} The repulsive core is the LJ form
with $z^{*}$ omitted,
\begin{equation}
  \Phi_{\mathrm{SDK}}(r) = \frac{48 d^{*12}}{r^{14}} - \frac{24 d^{*6}}{r^{8}},
  \qquad R_{\min}\le r \le d^{*}2^{1/6},
\end{equation}
augmented, in production runs only, by a cosine tail over
$d^{*}2^{1/6} < r \le R_{\mathrm{cut}}^{\mathrm{p}}$:
\begin{equation}
  \Phi_{\mathrm{SDK}} \mathrel{+}= \alpha_{\mathrm{SDK}}\,\varphi_{\mu\nu}
    \sin\!\left(\alpha_{\mathrm{SDK}} r^{2} + \beta_{\mathrm{SDK}}\right),
\end{equation}
\begin{equation}
  \alpha_{\mathrm{SDK}} = 1.530633312, \qquad \beta_{\mathrm{SDK}} = 1.213115524 ,
\end{equation}
where $\varphi_{\mu\nu}$ is a per-pair attractive strength supplied as an
$N\times N$ table. The \code{SDK\_stickers} variant is not implemented.

The equilibration flag is fixed for a run and baked into the
\code{ExcludedVolume} object at construction rather than passed per call.

% ===========================================================================
\chapter{Random number generation}
\label{ch:rng}

The generator is a bit-exact port of the Numerical Recipes \code{ran\_1} routine
used by the original Fortran, retained so that regression tests can compare the
two codes at machine precision. Each draw combines a Marsaglia xorshift with a
Park--Miller multiplicative congruential generator evaluated by Schrage's method.

With $\mathrm{IA}=16807$, $\mathrm{IM}=2^{31}-1$, $\mathrm{IQ}=127773$,
$\mathrm{IR}=2836$, the state $(i_x, i_y)$ evolves as
\begin{align}
  i_x &\leftarrow i_x \oplus (i_x \ll 13), \quad
  i_x \leftarrow i_x \oplus (i_x \gg 17), \quad
  i_x \leftarrow i_x \oplus (i_x \ll 5), \\
  k &= \lfloor i_y/\mathrm{IQ}\rfloor, \qquad
  i_y \leftarrow \mathrm{IA}\,(i_y - k\,\mathrm{IQ}) - \mathrm{IR}\,k,
  \qquad i_y \leftarrow i_y + \mathrm{IM}\;\text{if } i_y<0,
\end{align}
and the returned deviate is
\begin{equation}
  X = a_m \cdot \left[\left(\mathrm{IM} \;\&\; (i_x \oplus i_y)\right) \;|\; 1\right],
  \qquad
  a_m = \frac{\mathrm{nearest}(1,-1)}{\mathrm{IM}} .
\end{equation}
Seeding from $s$ sets $i_y = (888889999 \oplus |s|)\,|\,1$ and
$i_x = 777755555 \oplus |s|$.

Two implementation details matter for exactness. The right shift must be
\emph{logical} (zero-filling), matching Fortran's \code{ishft}, so $i_x$ is
manipulated as an unsigned 32-bit quantity. And $a_m$ is a \emph{single
precision} constant in the Fortran; keeping it single makes the stream
single-precision granular, and it is deliberately not promoted to double.

The generator state is encapsulated in an \code{Rng} object with
\code{save()}/\code{restore()}, rather than living in global module state. This
supports variance-reduction schemes, in which two chains must consume an
identical random stream, and makes segmented integration reproducible.

% ===========================================================================
\chapter{Initial configurations}
\label{ch:initial}

Initial configurations are generated in Python (\code{python/bdsim/initial.py}).
All samplers use inverse-transform sampling: for a one-dimensional density $p(X)$
with cumulative distribution $P$, if $Y\sim\mathcal{U}[0,1]$ then $X = P^{-1}(Y)$
is distributed as $p$.

\section{Gaussian chain}

Each bond is an isotropic Gaussian; bead positions are the cumulative sum,
recentred on the centre of mass. Adequate whenever the run equilibrates from the
starting state.

\section{FENE--Fraenkel equilibrium chain}

The equilibrium bond-length density for a FENE--Fraenkel connector follows from
the potential
\begin{equation}
  \phi_{\mathrm{eq}} = -\frac{H\,\delta Q^{2}}{2}
    \ln\!\left[1-\frac{(Q-\sigma)^{2}}{\delta Q^{2}}\right]
\end{equation}
via $\psi_{\mathrm{eq}}(Q) \propto Q^{2}e^{-\phi_{\mathrm{eq}}/\kB T}$, giving
\begin{equation}
  \psi_{\mathrm{eq}}(Q) \;\propto\; Q^{2}
    \left[1-\frac{(Q-\sigma)^{2}}{\alpha}\right]^{\alpha/2},
  \qquad \alpha = \delta Q^{2},
  \label{eq:ff-eq-dist}
\end{equation}
supported on $\max(0,\sigma-\delta Q) \le Q \le \sigma+\delta Q$. (The exponent is
$H\delta Q^{2}/2\kB T = \alpha/2$ in Hookean units.) The normalisation can be
obtained in closed form,
\begin{equation}
  J_{\mathrm{eq}} = \frac{\delta Q\left[\delta Q^{2}+(3+2c)\sigma^{2}\right]}{3+2c}\,
                    \mathrm{B}\!\left(\tfrac12,\,1+c\right),
  \qquad c = \frac{H\,\delta Q^{2}}{2\kB T},
\end{equation}
with $\mathrm{B}$ the Beta function, though the code does not need it: it
tabulates \eqref{eq:ff-eq-dist} on a grid of $10^{4}$ points, forms the cumulative
trapezoidal integral, normalises it to $[0,1]$, and inverts by linear
interpolation. Grid endpoints are inset by $10^{-7}$ to avoid the integrable
singularities in the derivative. Bond directions are then either isotropic on the
unit sphere (Marsaglia's rejection method) or all along $x$, giving
\code{fene\_fraenkel\_chain} and \code{fene\_fraenkel\_chain\_aligned\_x}
respectively.

\section{Correlated chains with a bending potential}
\label{sec:mc-chain}

For a chain with both a spring and a bending potential the two contributions to
the configurational integral separate, so bond lengths and bond angles may be
sampled independently and then assembled. Bond lengths come from
\eqref{eq:ff-eq-dist}; included angles come from \eqref{eq:angle-dist}, which for
the $C(1-\cos\theta)$ potential can be inverted analytically. With
$u = \cos\theta$ the density is $\propto e^{Cu}$ on $[-1,1]$, so
\begin{equation}
  u = \frac{1}{C}\ln\!\left[e^{-C} + Y\left(e^{C}-e^{-C}\right)\right],
  \qquad Y\sim\mathcal{U}[0,1],
  \label{eq:angle-sample}
\end{equation}
reducing to $u\sim\mathcal{U}[-1,1]$ when $C=0$.

The chain is then built sequentially. The first bond points isotropically. Each
subsequent bond is constructed in a local frame at polar angle $\theta_i$ from
the $z$-axis with a uniformly random azimuth, then rotated so that the local
$z$-axis aligns with the previous bond direction $\uv_{i-1}$; the rotation is the
Rodrigues matrix taking $\hat{\bm{z}}$ to $\uv_{i-1}$,
\begin{equation}
  \bm{R} = \bm{\delta} + [\bm{v}]_\times + [\bm{v}]_\times^{2}\frac{1-c}{\|\bm{v}\|^{2}},
  \qquad \bm{v} = \hat{\bm{z}}\times\uv_{i-1},
  \qquad c = \uv_{i-1}\cdot\hat{\bm{z}},
\end{equation}
with the degenerate cases $c\to\pm1$ handled separately. Multiplying by the
sampled length gives $\Q_i$, and bead positions follow by cumulative summation.
This reproduces bending correlations, unlike an isotropic random walk.


% ===========================================================================
\chapter{Coarse-graining a wormlike chain}
\label{ch:coarse}

The chapters so far describe a bead--spring chain with given parameters. This one
is the inverse problem: given a real semiflexible molecule --- a DNA fragment of
contour length $L$ and persistence length $l_p$ --- choose the spring and bending
parameters of an $N_s$-spring model that represents it. Everything here is
implemented in \code{python/bdsim/coarse\_grain.py}.

Units are unimportant as long as they are consistent: $L$ and $l_p$ may be in base
pairs, nanometres or metres, and only their ratio enters. DNA is quoted below as
$L = 25\,000$\,bp with $l_p = 147$\,bp (i.e.\ $50$\,nm at $0.34$\,nm per base pair).

\section{Wormlike-chain reference moments}
\label{sec:wlc-moments}

The target is the Kratky--Porod wormlike chain. Its even end-to-end moments follow
from the Hamprecht--Kleinert perturbation recursion for the rigid rotator in a
uniform field; the code carries them for $n = 1,\dots,6$ as reduced moments
\begin{equation}
  \langle r^{2n}\rangle = \frac{\langle R^{2n}\rangle}{L^{2n}},
  \qquad r = R/L,
  \qquad t = L/l_p .
\end{equation}
The lowest is the familiar
\begin{equation}
  \langle R^{2}\rangle = 2 l_p L - 2 l_p^{2}\left(1 - e^{-L/l_p}\right).
  \label{eq:wlc-r2}
\end{equation}
Each moment has the closed form
\begin{equation}
  \langle r^{2n}\rangle = t^{-2n}\sum_{l=0}^{n} P_{n,l}(t)\,e^{-l(l+1)t/2},
  \label{eq:wlc-closed}
\end{equation}
the exponential rates being the unperturbed rotator levels $e_0(l) = l(l+1)/2$.

\paragraph{Why the closed forms are not used everywhere.} For a stiff chain
($t\to0$) the $O(1)$ terms of \eqref{eq:wlc-closed} cancel to $O(t^{2n})$, so
evaluating it in double precision is catastrophically inaccurate: at $t = 10^{-3}$
the expression for $\langle R^{8}\rangle$ returns exactly $0$ instead of $\sim
10^{-24}$. The code therefore also stores the Taylor coefficients about $t=0$ and
switches between the two at $t = 0.8$, where they agree best. The worst-case
accuracy at the switch is $\sim10^{-13}$ for $n \le 3$, $\sim10^{-11}$ for $n=4$,
$\sim10^{-7}$ for $n=5$ and $\sim10^{-5}$ for $n=6$.

Both representations are verified against the two limits they must satisfy:
$\langle r^{2n}\rangle \to 1$ as $t\to0$ (a rigid rod has $R = L$), and the
Gaussian values $\langle R^{2n}\rangle \to \langle R^{2}\rangle^{n}(2n+1)!!/3^{n}$
as $t\to\infty$ --- e.g.\ $\langle R^{4}\rangle \to \tfrac{20}{3}(L l_p)^{2}$ and
$\langle R^{6}\rangle \to \tfrac{280}{9}(L l_p)^{3}$.

\section{Discretising the chain}

Each of the $N_s$ springs stands for a piece of contour of length
\begin{equation}
  l_s = \frac{L}{N_s},
  \qquad
  N_{\mathrm{K},s} = \frac{l_s}{2 l_p},
\end{equation}
$N_{\mathrm{K},s}$ being the number of Kuhn steps it contains. Two conditions fix
the spring.

\paragraph{(i) The spring is exactly as extensible as the contour it replaces.}
\begin{equation}
  \sigma + \delta Q = l_s .
  \label{eq:cg-extensibility}
\end{equation}

\paragraph{(ii) It has the right mean-square extension.}
\begin{equation}
  \langle Q^{2}\rangle_{\mathrm{FF}} = \langle R^{2}\rangle_{\mathrm{WLC}}(l_s, l_p),
  \label{eq:cg-secondmoment}
\end{equation}
with the left-hand side the second moment of the FENE--Fraenkel equilibrium
distribution \eqref{eq:ff-eq-dist}. Condition~\eqref{eq:cg-extensibility} removes
one parameter and \eqref{eq:cg-secondmoment} fixes a second, leaving the shape
exponent
\begin{equation}
  c = \frac{H \delta Q^{2}}{2 \kB T}
\end{equation}
to be supplied. The code solves \eqref{eq:cg-secondmoment} for $\sigma$ by
bisection with $c$ given, and falls back to $\sigma = 0$ (the FENE limit) when no
positive root exists.

\section{The shape exponent}
\label{sec:cg-exponent}

In the flexible limit Sunthar and Prakash give
\begin{equation}
  H = \frac{\kB T}{\delta Q^{2}}\left(\frac{3 l_s^{2}}{\langle R^{2}\rangle} - 5\right),
  \qquad\text{i.e.}\qquad
  c = \frac{1}{2}\left(\frac{3 l_s^{2}}{\langle R^{2}\rangle} - 5\right).
  \label{eq:cg-c-flexible}
\end{equation}
This becomes negative --- unphysical --- once
\begin{equation}
  \frac{3 l_s^{2}}{\langle R^{2}\rangle} < 5
  \quad\Longleftrightarrow\quad
  \frac{\langle R^{2}\rangle}{l_s^{2}} > \frac{3}{5}
  \quad\Longleftrightarrow\quad
  \frac{l_s}{l_p} < 1.772 ,
  \label{eq:cg-negative}
\end{equation}
i.e.\ for stiff segments.

A stiff segment must become \emph{harder}, not softer, so the correction has to
grow as $l_s/l_p \to 0$. Adding $l_p/l_s$ does that and vanishes in the flexible
limit, leaving \eqref{eq:cg-c-flexible} intact where it is valid:
\begin{equation}
  \boxed{\;
  c = \frac{1}{2}\left(\frac{3 l_s^{2}}{\langle R^{2}\rangle} - 5\right)
      + \frac{l_p}{l_s}\;}
  \label{eq:cg-c}
\end{equation}

The correction is chosen to act only where it is needed: $l_p/l_s$ diverges as the
segment becomes stiff and vanishes in the flexible limit, leaving
\eqref{eq:cg-c-flexible} untouched where it is valid. Without it $c$ would be
negative for every segment shorter than $1.772\,l_p$.

\section{Converting to simulation units}

The simulation works in Hookean units (Section~\ref{sec:units}), whose length scale
follows from the spring itself, $l_H = \sqrt{\kB T/H} = \delta Q/\sqrt{2c}$. Hence
\begin{equation}
  \boxed{\;
  \sqrt{b} \;=\; \delta Q_H = \sqrt{2c},
  \qquad
  \sigma_H = \frac{\sigma\sqrt{2c}}{\delta Q}\;}
  \label{eq:cg-hookean}
\end{equation}
which are exactly \code{SpringParams::sqrtb} and
\code{SpringParams::natural\_length}; the reduced natural length the solver uses is
$q_0 = \sigma_H/\sqrt{b}$. Note that $\delta Q_H^{2} = 2c = b$ recovers the FENE
$b$ parameter, as it must.

Stiff segments produce large $\sigma_H$ (e.g.\ $\sigma_H \approx 85$, $q_0\approx30$
for a $294$\,bp fragment with $N_s = 10$). These are numerically demanding springs
and need a correspondingly small timestep.

\section{Bending stiffness}
\label{sec:cg-bending}

The bending constant $C$ of \eqref{eq:bending-potential} controls the correlation
between successive springs. Four choices are implemented, and they answer
different questions:

\begin{itemize}
  \item \textbf{\code{saadat}} --- Saadat and Khomami's fit in $\nu = l_s/l_p =
        2N_{\mathrm{K},s}$,
        \begin{equation}
          C = \frac{1 + p_1\nu + p_2\nu^{2}}{\nu + p_3\nu^{2} + p_4\nu^{3}},
        \end{equation}
        with $(p_1,\dots,p_4) = (-1.237,\,0.8105,\,-1.0243,\,0.4595)$. Both
        numerator and denominator have negative discriminant, so $C>0$ for all
        $\nu>0$. This is the thesis choice.
  \item \textbf{\code{langevin}} --- match the continuous wormlike-chain tangent
        correlation exactly, $\langle\cos\theta\rangle = e^{-l_s/l_p}$, inverting
        the Langevin function \eqref{eq:angle-dist}.
  \item \textbf{\code{yamakawa}} --- match the \emph{discrete} chain's persistence
        length, $l_p/l_s = (1+\langle\cos\theta\rangle)/2(1-\langle\cos\theta\rangle)$.
  \item \textbf{\code{match\_chain}} --- solve for $C$ so that the assembled chain
        reproduces $\langle R^{2}\rangle_{\mathrm{WLC}}(L, l_p)$; see below.
\end{itemize}
The first three agree for $l_s \ll l_p$ and diverge as the segment becomes
flexible, because they match different things. In particular \code{saadat} is a fit
to the discretised chain's statistics, not to the local correlation, so its
$\langle\cos\theta\rangle$ is deliberately \emph{not} $e^{-l_s/l_p}$.

\section{Chain-level fidelity}
\label{sec:cg-chain}

Matching each segment to its own piece of contour does not make the assembled
chain reproduce $\langle R^{2}\rangle_{\mathrm{WLC}}(L,l_p)$, because the chain
value also depends on the angular correlation.

Write the end-to-end vector as the sum of the bonds, $\bm{R} = \sum_\mu\bm{Q}_\mu$,
and split each bond into its length and direction, $\bm{Q}_\mu = Q_\mu\uv_\mu$
with $Q_\mu = \|\bm{Q}_\mu\|$ a \emph{scalar}. Then
\begin{equation}
  \langle R^{2}\rangle = \sum_\mu \langle Q_\mu^{2}\rangle
    + 2\sum_{\mu<\nu}\langle \bm{Q}_\mu\cdot\bm{Q}_\nu\rangle .
\end{equation}
In this model the bond lengths are drawn independently of each other and of the
directions, so the cross terms factorise,
\begin{equation}
  \langle \bm{Q}_\mu\cdot\bm{Q}_\nu\rangle
    = \langle Q_\mu\rangle\,\langle Q_\nu\rangle\,\langle \uv_\mu\cdot\uv_\nu\rangle ,
\end{equation}
and the directions form a Markov chain: each step decorrelates by the same factor,
because the azimuthal angle about the previous bond is uniform, so
$\langle\uv_\mu\cdot\uv_\nu\rangle = c_\theta^{\,|\nu-\mu|}$ with
$c_\theta = \langle\cos\theta\rangle$. This is the classical freely-rotating-chain
argument, here carried through with a distributed rather than fixed bond length.
Summing the geometric series,
$\sum_{\mu<\nu}c_\theta^{\,\nu-\mu} = \sum_{k=1}^{N_s-1}(N_s-k)c_\theta^{k}
 = N_s c_\theta/(1-c_\theta) - c_\theta(1-c_\theta^{N_s})/(1-c_\theta)^{2}$,
gives
\begin{equation}
  \langle R^{2}\rangle_{\mathrm{chain}}
    = N_s\langle Q^{2}\rangle
    + 2\langle Q\rangle^{2}\left[
        \frac{N_s\,c_\theta}{1-c_\theta}
        - \frac{c_\theta\left(1-c_\theta^{N_s}\right)}{(1-c_\theta)^{2}}\right].
  \label{eq:cg-chain-r2}
\end{equation}
For a fixed bond length $b$ this reduces to the textbook freely-rotating-chain
result $\langle R^{2}\rangle = N_s b^{2}\left[(1+c_\theta)/(1-c_\theta)\right]
- 2b^{2}c_\theta(1-c_\theta^{N_s})/(1-c_\theta)^{2}$, which is a useful check on
the algebra.

\paragraph{A note on $\langle Q\rangle$.} Both $\langle Q^{2}\rangle$ and
$\langle Q\rangle$ here are moments of the scalar bond \emph{length}, not of the
bond vector --- the mean bond \emph{vector} is of course zero by isotropy. They are
not interchangeable: the diagonal terms carry the mean square length while the
cross terms carry the square of the mean length, and
\begin{equation}
  \langle Q^{2}\rangle - \langle Q\rangle^{2} = \operatorname{Var}(Q) > 0 ,
\end{equation}
so using $\langle Q^{2}\rangle$ in both places overestimates $\langle
R^{2}\rangle$ by exactly that variance in every cross term. Monte-Carlo sampling
confirms \eqref{eq:cg-chain-r2} is the closer of the two.

\paragraph{The size error this exposes.} With the \code{saadat} constant, the
assembled chain misses the wormlike-chain mean-square end-to-end distance by tens
of percent (Table~\ref{tab:cg-chain}) --- even though every segment is matched
exactly. \code{match\_chain} inverts \eqref{eq:cg-chain-r2} for $c_\theta$ (which
increases monotonically with it) and hence for $C$, bringing every case to within
about $1\%$. The trade-off is explicit: with this two-parameter family one can
have the correct local correlation or the correct global size, not both.

\begin{table}[ht]
\centering
\caption{Monte-Carlo $\langle R^{2}\rangle$ of the coarse-grained chain divided by
the wormlike-chain target, $l_p = 147$\,bp. Statistical error $\lesssim 0.01$.}
\label{tab:cg-chain}
\begin{tabular}{@{}rrrr@{}}
\toprule
$L$ (bp) & $N_s$ & \code{saadat} & \code{match\_chain} \\
\midrule
294    & 10 & 0.887 & 1.010 \\
294    & 30 & 0.953 & 1.000 \\
1223   & 20 & 0.722 & 1.003 \\
7003   & 30 & 1.102 & 1.002 \\
25000  & 30 & 1.023 & 0.996 \\
25000  & 80 & 1.152 & 0.993 \\
164000 & 40 & 0.996 & 0.996 \\
\bottomrule
\end{tabular}
\end{table}

\section{Which moments to fit}
\label{sec:cg-hk}

Matching $\langle Q^{2}\rangle$ is a choice, not a necessity. Hamprecht and
Kleinert give a criterion for which moments actually carry information: compare
each moment with what a \emph{flat} distribution on $r\in[0,1]$ would give,
$\langle r^{2n}\rangle^{\mathrm{flat}} = 1/(2n+2)$, and take the ratio
\begin{equation}
  \Xi(n) = \frac{\langle r^{2n}\rangle}{\langle r^{2n}\rangle^{\mathrm{flat}}}
         = (2n+2)\,\langle r^{2n}\rangle .
\end{equation}
$\Xi$ peaks at some order $n_{\max}$, roughly $4\xi$ with $\xi = l_p/L$, and those
are the informative moments: the low moments of a stiff chain all sit near $1$ and
say almost nothing about the shape. If $n_{\max}\le1$ the lowest moments
$\langle R^{2}\rangle$, $\langle R^{4}\rangle$, $\langle R^{6}\rangle$ are the
right ones.

Applied at chain level with $l_p = 147$\,bp, every DNA fragment from $239$\,bp
upwards has $n_{\max}\le1$, confirming that the lowest three moments are the
correct choice across the whole experimental range; only at $L\approx l_p$ does the
peak move (to $n=3$).

Applied at \emph{segment} level, where the relevant stiffness is $\xi_s = l_p/l_s$,
the picture changes, and \code{fit\_spring\_to\_moments} solves for
$(\sigma, c)$ --- two free parameters, hence exactly two moments --- to reproduce
the two most informative orders. Table~\ref{tab:cg-moments} shows that fitting the
informative moments rather than $\langle Q^{2}\rangle$ can change $c$ by an order
of magnitude, and does so precisely in the stiff regime where the correction of
Section~\ref{sec:cg-exponent} was doing the work. The method fails, honestly, when
the required order exceeds the tabulated $n\le6$ (very stiff segments); the moment
table can be regenerated to higher order from the recursion if needed.

\begin{table}[ht]
\centering
\caption{Segment-level two-moment fit against the default scheme, $l_p = 147$\,bp.
Residuals of the fit are $\sim10^{-15}$.}
\label{tab:cg-moments}
\begin{tabular}{@{}rrrrrrrr@{}}
\toprule
$L$ & $N_s$ & $l_s/l_p$ & peak $n$ & fitted orders & $\sigma/l_s$ & $c$ (fitted) & $c$ (default) \\
\midrule
294    & 10 & 0.20 & 6 & $(5,6)$ & 0.967 & 0.142 & 4.10 \\
1223   & 20 & 0.42 & 6 & $(5,6)$ & 0.933 & 0.515 & 1.62 \\
7003   & 30 & 1.59 & 2 & $(2,3)$ & 0.763 & 1.340 & 0.517 \\
25000  & 30 & 5.67 & 1 & $(1,2)$ & 0.333 & 3.874 & 2.835 \\
164000 & 40 & 27.9 & 1 & $(1,2)$ & 0.015 & 19.81 & 19.23 \\
\bottomrule
\end{tabular}
\end{table}

\section{From moments to a distribution}
\label{sec:cg-hkdist}

Hamprecht and Kleinert parameterise the end-to-end distribution as
\begin{equation}
  P(r) = \frac{\beta}{\mathrm{B}\!\left(\frac{3+k}{\beta},\,m+1\right)}\,
         r^{k+2}\left(1 - r^{\beta}\right)^{m},
  \qquad 0\le r\le 1,
  \label{eq:hk-pdf}
\end{equation}
a radial distribution (it already carries the $r^{2}$ Jacobian), whose moments are
closed-form,
\begin{equation}
  \langle r^{2n}\rangle =
    \frac{\Gamma\!\left(\frac{3+k+2n}{\beta}\right)\,
          \Gamma\!\left(\frac{3+k}{\beta}+m+1\right)}
         {\Gamma\!\left(\frac{3+k}{\beta}\right)\,
          \Gamma\!\left(\frac{3+k+2n}{\beta}+m+1\right)} .
  \label{eq:hk-moments}
\end{equation}
Fitting the three parameters $(k,\beta,m)$ to three moments therefore converts a
moment list into a closed-form distribution. The code does this by damped Newton
in $(k, \ln\beta, \ln(m+1))$, a parameterisation that enforces $\beta>0$ and
$m>-1$ automatically. The same machinery applies to the coarse-grained chain by
feeding it that chain's Monte-Carlo moments, so model and target can be compared
as smooth analytic curves rather than histograms. Fitted parameters agree with
independently tabulated values to four or five significant figures.

Two caveats. The fit is \emph{not unique}: at $L/l_p = 2$ two parameter sets were
found that both reproduce all three moments to $10^{-14}$ but differ elsewhere in
the distribution. And it fails for $L/l_p\lesssim1$, which is the paper's own
point --- there the informative moments are not the lowest three.


\section{Capping the spring constant}
\label{sec:cg-hcap}

Nothing above bounds $H$. As the discretisation is refined the segments become
stiff and $H$ grows without limit --- at $L = 2000$\,bp it reaches
$736$\,pN/nm by $N_s = 80$. That is not merely a large number: $H$ sets the time
unit through $\lambda_H = \zeta/4H$ (Section~\ref{sec:real-units}), so a timestep
of fixed size in simulation units covers proportionally \emph{fewer real seconds},
and reaching the chain's relaxation time becomes arbitrarily expensive.

The code therefore accepts a ceiling \code{max\_H}. When it binds, $c$ is pinned to
the cap, which through $\sigma+\delta Q = l_s$ ties $c$ to $\sigma$ by
$c = \tfrac12 H_{\max}(l_s-\sigma)^2$, leaving one unknown for the one remaining
condition --- so \eqref{eq:cg-secondmoment} is \emph{re-solved}, not abandoned.

The trade is cheap in exactly the regime where it bites. Once a segment is much
shorter than the persistence length its length barely fluctuates, and the chain's
conformation --- and hence its dynamics --- is governed by the bending potential
between segments rather than by the springs. Table~\ref{tab:hcap} bears this out:
capping $H$ by four orders of magnitude moves the mean segment length
$\sigma/l_s$ in the fourth decimal place, and $\langle Q^2\rangle$ still matches
the wormlike chain to machine precision. What changes is $\delta Q_H = \sqrt{2c}$,
the width of the length distribution within an unchanged physical range.

\begin{table}[ht]
\centering
\caption{Effect of capping $H$ at $0.05$\,pN/nm, for $L = 2000$\,bp,
$l_p = 147$\,bp. The second-moment condition is preserved throughout.}
\label{tab:hcap}
\begin{tabular}{@{}rrrrrrr@{}}
\toprule
 & & \multicolumn{2}{c}{uncapped} & \multicolumn{3}{c}{capped} \\
\cmidrule(lr){3-4}\cmidrule(lr){5-7}
$N_s$ & $l_s/l_p$ & $H$ (pN/nm) & $\sigma/l_s$ & $H$ (pN/nm) & $\sigma/l_s$ & $\langle Q^2\rangle$ error \\
\midrule
10 & 1.361 & 0.0175 & 0.7781 & --      & 0.7781 & $-1\times10^{-16}$ \\
20 & 0.680 & 0.499  & 0.8913 & 0.05    & 0.8872 & $+2\times10^{-16}$ \\
40 & 0.340 & 20.1   & 0.9453 & 0.05    & 0.9436 & $-4\times10^{-16}$ \\
80 & 0.170 & 736.5  & 0.9723 & 0.05    & 0.9718 & $-2\times10^{-16}$ \\
\bottomrule
\end{tabular}
\end{table}

% ===========================================================================
\chapter{Matching the dynamics}
\label{ch:dynamics}

The coarse-graining of Chapter~\ref{ch:coarse} fixes the chain's \emph{size} and
says nothing about how fast it moves. A model can therefore have the right
structure and completely the wrong timescale. This chapter closes that gap: it
chooses the bead hydrodynamic radius, through $h^{*}$, so the chain reproduces a
measured diffusivity --- and does so consistently across discretisations, which is
the difficult part. The implementation is \code{python/bdsim/dynamics.py}.

\section{The Kirkwood hydrodynamic radius}

In the code's units a free bead has $D = 1/4$ (the SDE carries $1/\sqrt2$ on the
noise, so $\langle\mathrm{d}R^2\rangle = 3\,\mathrm{d}t/2$), and the bead radius is
$a = \sqrt{\pi}h^{*}$. The Kirkwood approximation to the centre-of-mass
diffusivity is
\begin{equation}
  D_K = \frac{1}{4N^{2}}\left[N + \sum_{\mu\neq\nu}
        \left\langle \tfrac13\operatorname{tr}\bm{D}_{\mu\nu}\right\rangle\right],
  \label{eq:kirkwood-D}
\end{equation}
and the hydrodynamic radius, defined by $D = a/(4R_H)$ so that a single bead gives
$R_H = a$, is
\begin{equation}
  \boxed{\;
  R_H = \frac{a N^{2}}
             {N + \sum_{\mu\neq\nu}\left\langle\tfrac13\operatorname{tr}\bm{D}_{\mu\nu}\right\rangle}\;}
  \label{eq:kirkwood-RH}
\end{equation}

The trace collapses usefully. Substituting the RPY weights \eqref{eq:rpy-omegas}
into $\tfrac13\operatorname{tr}\bm{D}_{\mu\nu} = \omega_1 + \omega_2 r^{2}/3$, the
dyadic corrections cancel \emph{exactly} in both branches, leaving the Oseen form
\begin{equation}
  \tfrac13\operatorname{tr}\bm{D}_{\mu\nu} =
  \begin{cases}
    a/r, & r \ge 2a,\\[2pt]
    1 - r/4a, & r < 2a,
  \end{cases}
  \label{eq:rpy-trace}
\end{equation}
continuous at $r = 2a$ (both give $\tfrac12$). In the well-separated regime
\eqref{eq:kirkwood-RH} reduces to the familiar Kirkwood expression
\begin{equation}
  \frac{1}{R_H} = \frac{1}{Na} + \frac{1}{N^{2}}\sum_{\mu\neq\nu}
                  \left\langle\frac{1}{r_{\mu\nu}}\right\rangle,
  \label{eq:kirkwood-classic}
\end{equation}
which is linear in $1/a$ and can be inverted in closed form. The general case is
solved by bisection, since overlapping beads break that linearity.

\paragraph{Why this makes the problem tractable.} Equation~\eqref{eq:kirkwood-RH}
depends on the configurations only through $\langle\operatorname{tr}\rangle$ ---
that is, only on the \emph{static} structure, which Chapter~\ref{ch:coarse} has
already fixed. So $h^{*}$ follows from a Monte-Carlo average over equilibrium
configurations, with no dynamics run at all. It is a one-dimensional monotone
solve with a physical failure mode: as $h^{*}\to\infty$,
$R_H \to N^{2}/\sum\langle 1/r\rangle$, so a target above that non-draining
ceiling is unreachable. That is a real statement, not a numerical one --- it means
the discretisation is too coarse to carry the molecule's friction.

\section{Verification against direct simulation}

Equation~\eqref{eq:kirkwood-D} is exact for free draining ($D = 1/4N$, no
approximation) and approximate otherwise, so it must be checked. Measuring the
centre-of-mass diffusivity directly turned out to be a lesson in estimators.

A single-origin mean-square displacement, fitted over the later part of each run,
gave $18\%$ relative error --- useless for testing a few-percent approximation.
Accumulating over \emph{all} time origins helped less than expected ($12\%$),
because displacements at lag $\tau$ overlap: a run of length $T$ contains only
about $T/\tau$ independent samples however many origins are used. The fix is to
use \emph{short} lags, which is also the right definition: the Kirkwood formula is
an equilibrium-averaged instantaneous mobility, i.e.\ the short-time diffusivity,
and the long-time value is slightly lower. With that, Table~\ref{tab:kirkwood-bd}.

\begin{table}[ht]
\centering
\caption{Kirkwood estimate against direct Brownian dynamics, $N = 10$ Hookean
chain. The $h^{*} = 0$ row is exact and so tests the estimator rather than the
physics.}
\label{tab:kirkwood-bd}
\begin{tabular}{@{}rrrr@{}}
\toprule
$h^{*}$ & $D_K$ & $D$ measured & ratio \\
\midrule
0.00 & 0.025000 & $0.025500 \pm 0.00089$ & 1.020 \\
0.15 & 0.054322 & $0.054395 \pm 0.00191$ & 1.001 \\
0.30 & 0.081580 & $0.081088 \pm 0.00284$ & 0.994 \\
\bottomrule
\end{tabular}
\end{table}

\section{Real units}
\label{sec:real-units}

Matching a diffusivity quoted in simulation units would be meaningless, because
\emph{both} Hookean units move with the coarse-graining:
\begin{equation}
  l_H = \sqrt{\kB T/H}\quad\text{(set by the static fit)},
  \qquad
  \lambda_H = \frac{\zeta}{4H},\quad \zeta = 6\pi\eta_s a,\quad a = \sqrt{\pi}h^{*}l_H .
\end{equation}
$\lambda_H$ depends on $h^{*}$, so each discretisation counts time in a different
second. Across $N_s = 5$ to $20$ for a $2$\,kbp fragment $\lambda_H$ spans
$3.3\times10^{-5}$ to $5.3\times10^{-8}$\,s.

What makes the match well posed is that $R_H$ is a \emph{length}. The chain of
reasoning is:
\begin{enumerate}
  \item the static fit fixes $l_H$ in nanometres, $l_H = \delta Q/\sqrt{2c}$;
  \item the model's $R_H$ therefore converts to nanometres with no reference to
        time at all, and this is what is matched;
  \item Stokes--Einstein then gives the real diffusivity,
        \begin{equation}
          D = \frac{\kB T}{6\pi\eta_s R_H},
          \label{eq:stokes-einstein}
        \end{equation}
        so $D$ is matched in $\mathrm{m^2/s}$ the moment $R_H$ is matched in
        metres, independently of $N_s$;
  \item $\lambda_H$ in seconds then follows, and converts any simulated time into
        a real one.
\end{enumerate}
Substituting $a = \sqrt{\pi}h^{*}l_H$ and $R_H = a/4D_H$ into
$D = D_H l_H^{2}/\lambda_H$ reduces it algebraically to
\eqref{eq:stokes-einstein}; numerically the two routes agree to a ratio of
$1.0000$, which is the check that the unit bookkeeping is right.

The three conversions needed to interpret any output are
\begin{equation}
  l_H \;\text{(length)},
  \qquad
  F_H = \sqrt{H\kB T} = \frac{\kB T}{l_H} \;\text{(force)},
  \qquad
  \lambda_H \;\text{(time)},
\end{equation}
with energies in $\kB T$. Table~\ref{tab:real-units} shows them for a $2$\,kbp fragment.

It is worth being explicit about what is invariant there, since no two of those
columns cancel on their own. The physical diffusivity is
\begin{equation}
  D = D_K\,\frac{l_H^{2}}{\lambda_H},
\end{equation}
and across the three rows the three factors change by
\begin{equation}
  D_K \;\div\; 2.6,
  \qquad
  l_H^{2} \;\div\; 268,
  \qquad
  \lambda_H \;\div\; 689,
\end{equation}
so that the \emph{combination} $D_K l_H^{2}/\lambda_H$ is unchanged to four
figures --- $2.6\times268/689 = 1.00$. It has to be: that combination equals
$\kB T/6\pi\eta_s R_H$ by \eqref{eq:stokes-einstein}, and $R_H$ was matched.

\begin{table}[ht]
\centering
\caption{Unit conversions and the resulting physical diffusivity, $L = 2000$\,bp,
target $R_H = 42.7$\,nm.}
\label{tab:real-units}
\begin{tabular}{@{}rrrrrrr@{}}
\toprule
$N_s$ & $h^{*}$ & $l_H$ (nm) & $F_H$ (pN) & $\lambda_H$ (s) & $D_K$ (code) & $D$ ($\mu$m$^2$/s) \\
\midrule
5  & 0.1234 & 47.04 & 0.0875 & $2.61\times10^{-5}$ & 0.0603 & 5.1179 \\
10 & 0.2377 & 15.32 & 0.2686 & $1.74\times10^{-6}$ & 0.0378 & 5.1179 \\
20 & 0.7875 & 2.871 & 1.434  & $3.78\times10^{-8}$ & 0.0235 & 5.1179 \\
\bottomrule
\end{tabular}
\end{table}



\section{Free draining: matching without hydrodynamic interaction}
\label{sec:free-draining}

Hydrodynamic interaction dominates the cost of a step --- Table~\ref{tab:hi-cost}
puts it at forty to sixty times everything else --- so it is useful to be able to
switch it off for exploratory runs. Doing so changes the matching problem
completely, and for the better.

With $h^{*} = 0$ the diffusion tensor is the identity, so
\begin{equation}
  D = \frac{\kB T}{N\zeta} = \frac{\kB T}{6\pi\eta_s N a},
\end{equation}
independent of configuration. Matching a measured diffusivity therefore needs no
Monte-Carlo sampling and no root find; the bead radius follows in closed form,
\begin{equation}
  \boxed{\; a = \frac{R_H}{N} \;}
  \label{eq:fd-a}
\end{equation}
since $R_H = \kB T/6\pi\eta_s D$. Equivalently $R_H = Na$, which is the
$h^{*}\to0$ limit of \eqref{eq:kirkwood-RH}: with no shielding, the chain's
friction is simply $N$ times a bead's.

\paragraph{A nominal $h^{*}$.} The simulation is run with $h^{*} = 0$ and never
builds the RPY tensor. But the bead radius still has to be declared, because it
fixes the friction and hence the time unit $\lambda_H = \zeta/4H$. So
\eqref{eq:fd-a} yields a \emph{nominal} $h^{*}$, used only for the unit
conversions and never passed to the integrator. Running with $h^{*} = 0$ while
carrying $h^{*}_{\mathrm{nominal}}$ in the bookkeeping is the whole of the trick.

\begin{table}[ht]
\centering
\caption{Free-draining matching for $\lambda$-DNA, target $R_H = 328.2$\,nm. The
diffusivity is exact for every discretisation, with no sampling.}
\label{tab:free-draining}
\begin{tabular}{@{}rrrrrrr@{}}
\toprule
$N_s$ & $a$ (nm) & $h^{*}_{\mathrm{nom}}$ & $l_H$ (nm) & $\lambda_H$ (s) & $D$ ($\mu$m$^2$/s) & $\tau_1$ (s) \\
\midrule
10 & 29.84 & 0.0693 & 242.7 & $2.01\times10^{-3}$ & 0.6654 & 0.0994 \\
20 & 15.63 & 0.0498 & 177.2 & $5.62\times10^{-4}$ & 0.6654 & 0.1006 \\
30 & 10.59 & 0.0402 & 148.4 & $2.67\times10^{-4}$ & 0.6654 & 0.1041 \\
40 & 8.005 & 0.0346 & 130.4 & $1.56\times10^{-4}$ & 0.6654 & 0.1062 \\
\bottomrule
\end{tabular}
\end{table}

\paragraph{What is lost.} Only the centre-of-mass motion is constrained by
\eqref{eq:fd-a}. The internal modes of a free-draining chain follow Rouse rather
than Zimm dynamics and couple to flow differently, so the rheology is not to be
trusted.

The longest relaxation time is nonetheless a useful number, and for a free-draining
Hookean chain it needs no fitted prefactor,
\begin{equation}
  \tau_1 = \frac{\lambda_H}{\sin^{2}\!\left(\pi/2N\right)} .
\end{equation}
Table~\ref{tab:free-draining} shows it coming out close to
discretisation-independent, and in the right range for $\lambda$-DNA. The obvious
worry --- that Rouse $N^{2}$ scaling would inflate it --- does not materialise,
because $\lambda_H$ shrinks in step ($a = R_H/N$, and $l_H$ falls as the chain is
refined) and the two nearly cancel. That is a coincidence about $\tau_1$ alone,
not a reprieve for the rest of the spectrum: use free draining to place a run at a
sensible Weissenberg number and to shake down a workflow, and measure rheology
with HI.

\section{Rodlike units}
\label{sec:rodlike}

Hookean units are built from the spring, $l_H = \sqrt{\kB T/H}$. That is awkward
for a stiff, nearly rigid spring: $H$ is then large, $l_H$ is far smaller than the
bond it describes, and every length in the problem comes out as a large number
($\sigma_H \approx 35$ for a $2$\,kbp fragment at $N_s = 40$). It also makes
$h^{*}$ hard to read, for the reason given above.

The alternative is to scale lengths by something taken from the chain rather than
from the spring. Two choices are natural, and the code supports both:
\begin{equation}
  l_R = \sigma \quad\text{(the natural length)},
  \qquad\text{or}\qquad
  l_R = \sqrt{\langle Q^{2}\rangle} \quad\text{(the r.m.s.\ bond length)}.
\end{equation}
The first is undefined for a FENE spring, where $\sigma = 0$; the second is always
available, and the two coincide as the spring becomes rigid. Time and force follow
in the usual diffusive way,
\begin{equation}
  \lambda_R = \frac{\zeta\, l_R^{2}}{\kB T},
  \qquad
  F_R = \frac{\kB T}{l_R},
\end{equation}
matching the convention $\lambda^{*}_R = L^{2}\zeta/\kB T$ used for rodlike models,
with $L$ the rod length.

\subsection{Conversion}

Let $s = l_R/l_H$ be the rodlike length unit expressed in Hookean units --- that
is, $s = \sigma_H$ or $s = \sqrt{\langle Q^{2}\rangle_H}$, both of which the code
already has. Since $H = 1$ and $\kB T = 1$ in Hookean units,
\begin{equation}
  \frac{\lambda_H}{\lambda_R}
    = \frac{\zeta/4H}{\zeta l_R^{2}/\kB T}
    = \frac{\kB T}{4 H l_R^{2}}
    = \frac{l_H^{2}}{4 l_R^{2}}
    = \frac{1}{4s^{2}},
\end{equation}
the factor $4$ being the one in $\lambda_H = \zeta/4H$. Collecting the three
scales:
\begin{equation}
  \boxed{\;
  X_R = \frac{X_H}{s} \;\text{(length)},
  \qquad
  F_R = s\,F_H \;\text{(force)},
  \qquad
  t_R = \frac{t_H}{4s^{2}} \;\text{(time)},
  \qquad
  \dot\gamma_R = 4s^{2}\,\dot\gamma_H \;}
\end{equation}
with energies still in $\kB T$ and the bending constant $C$ already dimensionless
and therefore unchanged.

\subsection{Parameters}

The spring parameters transform to
\begin{equation}
  H^{*}_R = \frac{H l_R^{2}}{\kB T} = s^{2},
  \qquad
  \sigma^{*}_R = \frac{\sigma_H}{s},
  \qquad
  \delta Q^{*}_R = \frac{\delta Q_H}{s},
\end{equation}
and the hydrodynamic parameter, using $a = \sqrt{\pi}h^{*}l_H$ and the rodlike
convention $h^{*}_R = 3a/4l_R$, to
\begin{equation}
  h^{*}_R = \frac{3\sqrt{\pi}\,h^{*}}{4 s}.
\end{equation}

Scaling by $\sigma$ gives three identities worth using as checks, all of which the
implementation satisfies exactly:
\begin{equation}
  \sigma^{*}_R = 1,
  \qquad
  H^{*}_R = \sigma_H^{2},
  \qquad
  \delta Q^{*}_R = \frac{\delta Q_H}{\sigma_H} = \frac{1}{q_0},
\end{equation}
the last being the inverse of the reduced natural length the implicit solver works
with (Chapter~\ref{ch:springs}). For a $2$\,kbp fragment at $N_s = 40$
($\sigma_H = 35.53$, $\delta Q_H = 2.057$, $q_0 = 17.28$) the two scalings give
$H^{*}_R = 1262.7$ versus $1265.6$ and $\delta Q^{*}_R = 0.05788$ versus
$0.05781$ --- indistinguishable, as they must be once the spring is stiff enough
that $\sigma \simeq \sqrt{\langle Q^{2}\rangle}$.

\section{Experimental calibration for DNA}

The reference scale is Zimm's result for an ideal non-draining coil,
\begin{equation}
  D_{\mathrm{Zimm}} = \frac{8}{3\sqrt{6\pi^{3}}}\,
                      \frac{\kB T}{\eta_s\sqrt{L b}},
  \qquad b = 2 l_p,
  \label{eq:zimm}
\end{equation}
the prefactor being $0.195510$. Real DNA departs from \eqref{eq:zimm} in both
directions: short fragments are rod-like and diffuse faster, while long ones are
swollen by excluded volume and diffuse slower, so $D/D_{\mathrm{Zimm}}$ falls
below unity and keeps drifting as $L^{1/2-\nu}$ with $\nu\approx0.588$. The code
tabulates that ratio and interpolates it logarithmically, giving
\begin{equation}
  D(L) = D_{\mathrm{Zimm}}(L)\times\left[D/D_{\mathrm{Zimm}}\right](L),
  \qquad
  R_H = \frac{\kB T}{6\pi\eta_s D}.
\end{equation}

\paragraph{Provenance of the table.} The tabulated ratios were read off a
published figure (dsDNA with excluded volume, fitted by a discrete wormlike chain
with $b/d = 36$, $d = 2.9$\,nm, against dynamic light scattering, sedimentation
and single-molecule data). They are good to perhaps a few percent and exist so the
pipeline runs end to end; every entry point accepts a replacement table. Note also
that the digitised curve is the one fitting the light-scattering and sedimentation
data --- single-molecule values sit systematically below it, which the source
attributes to intercalating dyes. For $\lambda$-DNA this route gives
$0.665\,\mu$m$^2$/s where single-molecule measurements give $\approx0.47$; the
difference is the choice of data, not an error.

\begin{table}[ht]
\centering
\caption{The calibration across the experimental range, $l_p = 147$\,bp,
$b = 100$\,nm, water at $298$\,K.}
\label{tab:dna-calib}
\begin{tabular}{@{}rrrrrr@{}}
\toprule
$L$ (bp) & $L$ ($\mu$m) & $D_{\mathrm{Zimm}}$ & $D/D_{\mathrm{Zimm}}$ & $D$ ($\mu$m$^2$/s) & $R_H$ (nm) \\
\midrule
2\,000   & 0.68  & 3.086 & 1.658 & 5.118 & 42.7 \\
10\,000  & 3.40  & 1.380 & 1.293 & 1.784 & 122.4 \\
48\,502  & 16.5  & 0.627 & 1.062 & 0.665 & 328 \\
164\,000 & 55.8  & 0.341 & 0.954 & 0.325 & 672 \\
\bottomrule
\end{tabular}
\end{table}

\section{What is and is not discretisation-independent}

$h^{*}$ varies strongly with $N_s$ --- from $0.12$ to $0.79$ in
Table~\ref{tab:real-units} --- and must, since a chain of few large beads shields
itself differently from one of many small ones. The physical prediction is what
should not vary, and $D$ does not.

\paragraph{Large $h^{*}$ is not by itself a warning.} The familiar rule
$h^{*}\lesssim0.3$ assumes flexible springs, for which $l_H$ is comparable to the
bond length. Stiff FENE--Fraenkel springs have $l_H \ll \langle Q\rangle$, so
$h^{*}>1$ can be perfectly physical. The meaningful diagnostic is $a/\langle
Q\rangle$, the bead radius relative to the bond length; beyond about $0.5$ the
beads overlap and the RPY tensor is being used outside its intended regime. For
the cases above it stays between $0.14$ and $0.24$.

\paragraph{A residual drift.} $R_H$ is matched by construction and
$\langle R^{2}\rangle$ by \code{match\_chain}, but $\langle R_g^{2}\rangle$ is
not, and it drifts: $R_g/R_H$ falls from $2.10$ to $1.93$ as $N_s$ goes from $5$
to $40$. Matching one static moment does not fix the shape, so shape-sensitive
observables retain a discretisation dependence at the ten-percent level. Whether
the longest relaxation time inherits that drift is the genuine test of this whole
scheme --- the diffusivity is matched by construction and so proves nothing --- and
is what \code{validation/validate\_dynamics.py --relax} is for.

% ===========================================================================
\chapter{Observables}
\label{ch:observables}

All observables are computed in Python from a configuration $\R$ and, where
required, the bead forces $\bm{F}$. Let $\bm{r}_c = \frac1N\sum_\nu\bm{r}_\nu$ and
$\bm{r}'_\nu = \bm{r}_\nu-\bm{r}_c$.

\paragraph{Size measures.}
\begin{equation}
  R^{2} = \left\|\bm{r}_N-\bm{r}_1\right\|^{2},
  \qquad
  \bm{G} = \frac{1}{N}\sum_\nu \bm{r}'_\nu\bm{r}'_\nu,
  \qquad
  R_g^{2} = \operatorname{tr}\bm{G} .
\end{equation}
The $x$-stretch reported in comparisons is
$\max_\nu r_{\nu,x} - \min_\nu r_{\nu,x}$.

\paragraph{Kramers--Kirkwood stress.} The polymer contribution to the stress
tensor for one chain is
\begin{equation}
  \tau_{ij} = \sum_{\nu} r'_{\nu,i}\,F_{\nu,j},
  \label{eq:kramers}
\end{equation}
computed as \code{Rc.T @ F}. Which force is passed determines what the stress is
attributed to --- the spring force alone, or the total intramolecular force. From
\eqref{eq:kramers},
\begin{equation}
  \eta_p = -\frac{\tau_{xy}}{\dot\gamma},
  \qquad
  N_1 = \tau_{xx}-\tau_{yy},
  \qquad
  N_2 = \tau_{yy}-\tau_{zz}.
\end{equation}

\paragraph{The viscometric relaxation time.} The zero-shear polymer viscosity
defines a relaxation time,
\begin{equation}
  \lambda_\eta = \frac{M\,\eta_{p,0}}{c\,N_A\,\kB T} = \frac{\eta_{p,0}}{n\,\kB T},
  \label{eq:lambda-eta}
\end{equation}
with $M$ the molar mass, $c$ the mass concentration and $n = cN_A/M$ the chain
number density. The concentration cancels, so this is a single-chain property. It
is exactly what the code reports: the Kramers stress is measured in units of
$\kB T$, so $-\langle\tau_{xy}\rangle/\dot\gamma$ is already
$\eta_p/(n\kB T)$, a time in simulation units. In terms of the intrinsic viscosity
$[\eta] = \eta_{p,0}/(c\eta_s)$, equation~\eqref{eq:lambda-eta} becomes
$\lambda_\eta = M\eta_s[\eta]/(N_A\kB T)$.

\paragraph{Rouse check, and a factor of two.} $\lambda_\eta$ sums the whole
relaxation spectrum, not the slowest mode. For a free-draining Hookean chain the
mode amplitudes decay with $\tau_p = 1/\sin^{2}(p\pi/2N)$ --- this is what the
end-to-end correlation measures --- but the stress is \emph{quadratic} in the
amplitudes, so stress correlations decay at twice the rate, with time constant
$\tau_p/2$. Hence
\begin{equation}
  \lambda_\eta = \frac12\sum_{p=1}^{N-1}\frac{1}{\sin^{2}(p\pi/2N)}
               = \frac{N^{2}-1}{3},
  \label{eq:rouse-eta}
\end{equation}
which tends to $(\pi^{2}/12)\,\tau_1 \approx 0.82\,\tau_1$: \emph{shorter} than the
slowest mode. Omitting the factor of one half gives $2(N^2-1)/3$ and is wrong by
exactly two. A direct measurement of a Hookean chain at $N=11$ gives
$42.5 \pm 2.6$, against $40.0$ from \eqref{eq:rouse-eta} and $80.0$ without the
factor. The equilibrium validation script uses \eqref{eq:rouse-eta} as an analytic
target.


% ===========================================================================
\chapter{Force-extension}
\label{ch:force-extension}

\section{Measuring it, rather than inferring it}

Given the fit of Section~\ref{sec:cg-hkdist} one is tempted to read off a
force-extension curve. The free energy of a chain whose ends are held a distance
$R$ apart is $A(R) = -\kB T\ln[P(R)/R^{2}]$ (dividing out the angular Jacobian), so
$f = \mathrm{d}A/\mathrm{d}R$ gives
\begin{equation}
  \frac{f(r) L}{\kB T} = \frac{m\beta r^{\beta-1}}{1-r^{\beta}} - \frac{k}{r} .
  \label{eq:hk-force}
\end{equation}
This has the right Gaussian limit, $f = 3\kB T R/\langle R^{2}\rangle$ when
$k\to0$, $\beta\to2$.

\paragraph{It should not be trusted quantitatively.} Equation~\eqref{eq:hk-force}
is a \emph{consequence} of a three-moment fit, and low moments constrain the bulk
of the distribution, not the $r\to1$ tail that sets the high-force response. Its
divergence as $r\to1$ has the right form only because that form was built into the
ansatz \eqref{eq:hk-pdf}, with $m$ fitted to moments that barely see the region.
The $-k/r$ term likewise makes it meaningless as $r\to0$. It is useful for
comparing two distributions on equal footing, not for predicting a stretching
curve.

The honest alternative is to apply the forces and measure the response, which is
what the external force of Section~\ref{sec:external} is for: equal and opposite
loads \eqref{eq:stretch} on the two ends, equilibrate, and sample $\langle
R_z\rangle$. That measurement includes the bending potential, finite
extensibility, chain discreteness and hydrodynamic interaction, none of which
\eqref{eq:hk-force} knows about. It is implemented in
\code{validation/force\_extension.py}.

\paragraph{Equilibration is the trap.} A pulling measurement must be equilibrated
for several times the longest relaxation time,
\begin{equation}
  \lambda_1 = \frac{1}{\sin^{2}\!\left(\pi/2N\right)}
\end{equation}
in the code's time unit (the free-draining Hookean value; stiffer springs relax
faster, so this is a conservative bound). For $N = 11$ this is $\lambda_1\approx49$
--- and equilibrating for, say, $t = 10$ produces extensions three to ten times
too small, which looks exactly like a model that is far too stiff. The script
computes $\lambda_1$ and defaults to $t_{\mathrm{eq}} = 3\lambda_1$,
$t_{\mathrm{run}} = 6\lambda_1$.

\paragraph{An exact check.} For a Hookean chain every spring carries the full
applied load, so $\langle R_z\rangle = N_s f$ exactly, independent of $N_s$.
Measured with $N_s = 10$: ratios to the exact value of $0.956$, $0.975$, $0.984$,
$0.989$ at $f = 0.25,\,0.5,\,1,\,2$, all within one standard error.

\section{Does the model follow Marko--Siggia?}
\label{sec:ms}

At low force, yes, and necessarily so: the Marko--Siggia interpolation
\begin{equation}
  \frac{f l_p}{\kB T} = \frac{1}{4(1-x)^{2}} - \frac{1}{4} + x,
  \qquad x = R/L,
  \label{eq:marko-siggia}
\end{equation}
expands to $f l_p/\kB T \simeq 3x/2$, i.e.\ $f = 3\kB T R/\langle
R^{2}\rangle$ --- the Gaussian result. Any model with the correct $\langle
R^{2}\rangle$, which \code{match\_chain} enforces, reproduces it.

At high extension, no --- for two independent reasons.

\paragraph{(i) The spring law has the wrong divergence.} The FENE--Fraenkel force
diverges linearly in the gap to full extension,
$\bm{F}^{\mathrm{c}}\propto\left[\delta Q - (Q-\sigma)\right]^{-1}$, giving
$1-x\propto f^{-1}$. The wormlike chain's $1/4(1-x)^{2}$ gives $1-x\propto
f^{-1/2}$. This is intrinsic to FENE-type springs and no choice of parameters
fixes it.

\paragraph{(ii) Discretisation truncates the fluctuation spectrum.} The
$f^{-1/2}$ law comes from equipartition over transverse bending modes down to
arbitrarily short wavelength. A chain of $N_s$ springs has no modes below $l_s$, so
once $\sqrt{\kappa/f} < l_s$ the model has run out of fluctuations to pull out. The
crossover is
\begin{equation}
  \frac{f^{*} l_p}{\kB T} \sim \left(\frac{l_p}{l_s}\right)^{2},
  \label{eq:fstar}
\end{equation}
which for $L = 25\,000$\,bp with $N_s = 30$ is only $\approx0.03$: the model
departs from \eqref{eq:marko-siggia} almost immediately.

Table~\ref{tab:fe} shows both the measured curve and an independent
semi-analytic calculation (the single spring in the fixed-force ensemble, which
omits bending). They agree on the essential number and both differ from
Marko--Siggia in the same direction: the model over-extends, its extension deficit
falling a decade per decade of force ($f^{-1}$) against Marko--Siggia's factor of
$3.16$ ($f^{-1/2}$).

\begin{table}[ht]
\centering
\caption{Extension deficit $1-x$ for a coarse-grained $25\,000$\,bp fragment.
``Measured'' is a pulled Brownian-dynamics run including bending; ``calculated''
is the single-spring fixed-force ensemble without bending.}
\label{tab:fe}
\begin{tabular}{@{}rrrr@{}}
\toprule
$f l_p/\kB T$ & measured & calculated & Marko--Siggia \\
\midrule
1   & $5.2\times10^{-1}$ & $5.5\times10^{-1}$ & $5.6\times10^{-1}$ \\
10  & $7.6\times10^{-2}$ & $8.2\times10^{-2}$ & $1.6\times10^{-1}$ \\
100 & $7.8\times10^{-3}$ & $8.5\times10^{-3}$ & $5.0\times10^{-2}$ \\
\bottomrule
\end{tabular}
\end{table}

\paragraph{What to do if the high-force regime matters.} Use a spring law that has
the right divergence: the code carries Marko--Siggia (\code{WLC}) and bounded-WLC
(\code{WLCbounded}) force laws (Table~\ref{tab:force-laws}). The parameters cannot
simply be carried over, though --- the correct coarse-grained spring law is not the
underlying wormlike-chain law, and depends on $N_{\mathrm{K},s}$. Alternatively,
increase $N_s$ until $l_s\sim l_p$, which pushes $f^{*}$ up by
\eqref{eq:fstar}. For flows that never reach $f^{*}$, the wrong asymptote is
irrelevant and FENE--Fraenkel is entirely adequate.


% ===========================================================================
\chapter{Error bars for correlated data}
\label{ch:statistics}

Every number this code produces is a time average over a steady-state trajectory,
and successive samples of such a trajectory are not independent: a chain remembers
its configuration for roughly its longest relaxation time. Treating $N$ samples as
$N$ independent measurements understates the error by $\sqrt{g}$, where $g$ is the
statistical inefficiency --- often a factor of several, and an order of magnitude
in bad cases. The tools are in \code{python/bdsim/statistics.py}.

\section{Statistical inefficiency}

Define the normalised autocorrelation function $\rho_k$ of the sampled series. The
variance of the mean is inflated by
\begin{equation}
  g = 1 + 2\sum_{k\ge1}\rho_k,
  \qquad
  N_{\mathrm{eff}} = \frac{N}{g},
  \qquad
  \mathrm{s.e.} = \frac{\mathrm{std}(x)}{\sqrt{N_{\mathrm{eff}}}} ,
  \label{eq:stat-ineff}
\end{equation}
$g$ being twice the integrated autocorrelation time in units of the sampling
interval. The sum must be truncated: the tail of $\rho_k$ is pure noise and
summing it produces a divergent estimate. The code uses Sokal's automatic window
--- take the smallest $W$ with $W \ge c\,\tau_{\mathrm{int}}(W)$, here $c = 6$.

The simple protocol of sampling once per relaxation time and treating the samples
as independent is the $g \to$ (sampling interval) limit of the same idea. It is
correct but wasteful: \eqref{eq:stat-ineff} recovers the information in the
intervening samples instead.

\section{Blocking}

Flyvbjerg--Petersen blocking provides an independent estimate that assumes nothing
about the shape of $\rho_k$. Average adjacent pairs repeatedly; the apparent
standard error rises and then plateaus once blocks exceed the correlation time,
and the plateau is the answer. The code reports it alongside
\eqref{eq:stat-ineff} precisely so the two can be compared.

\section{Discarding the transient}

How much of the start to throw away is chosen by maximising the effective sample
size of what remains (Chodera's rule), trading lost samples against reduced bias.
This is not cosmetic: retaining a slowly-relaxing transient biases the mean
\emph{and} inflates the apparent correlation time, so the error bar grows too.

\section{Validation}

The estimators are checked against a process with an exact answer rather than
against each other. A first-order autoregressive series, $x_{t+1} = \phi x_t +
\varepsilon_t$, has $\rho_k = \phi^{k}$ and therefore
\begin{equation}
  g = \frac{1+\phi}{1-\phi}
\end{equation}
exactly. Table~\ref{tab:ar1} compares the estimate with that value, and the
predicted standard error with the measured scatter of the mean over $400$
independent realisations.

\begin{table}[ht]
\centering
\caption{Estimators against the exact AR(1) result. ``true'' is the empirical
scatter of the mean over 400 independent runs; ``naive'' ignores correlation.}
\label{tab:ar1}
\begin{tabular}{@{}rrrrrrr@{}}
\toprule
$\phi$ & $g$ exact & $g$ est. & ratio & s.e.\ true & s.e.\ from $g$ & s.e.\ naive \\
\midrule
0.00 & 1.00  & 1.01   & 1.011 & 0.0165 & 0.0169 & 0.0162 \\
0.50 & 3.00  & 3.01   & 1.004 & 0.0320 & 0.0325 & 0.0183 \\
0.80 & 9.00  & 9.05   & 1.005 & 0.0775 & 0.0890 & 0.0259 \\
0.90 & 19.00 & 19.28  & 1.015 & 0.1480 & 0.1454 & 0.0342 \\
0.95 & 39.00 & 37.21  & 0.954 & 0.3147 & 0.2591 & 0.0463 \\
0.98 & 99.00 & 102.88 & 1.039 & 0.8151 & 0.8500 & 0.0775 \\
\bottomrule
\end{tabular}
\end{table}

$g$ is recovered to within a few percent over two orders of magnitude, and the
naive error is wrong by $\sqrt{g}$ --- a factor of ten at $\phi = 0.98$.

\section{Two levels of averaging, and when to distrust the correction}

An ensemble run has correlated samples \emph{within} each trajectory and genuinely
independent trajectories. Two estimates are therefore available: the
autocorrelation-corrected pooled error \eqref{eq:stat-ineff}, and the scatter of
the per-trajectory means. The second needs no assumptions at all but has only
$n_{\mathrm{traj}}-1$ degrees of freedom; agreement between them is good evidence
that both can be trusted.

\paragraph{A silent failure mode.} The autocorrelation estimate is only meaningful
when each trajectory contains many samples per correlation time. With few samples
the ACF cannot be measured at all: Sokal's window closes on the first step, $g$
comes back as $\approx 1$ --- ``uncorrelated'' --- and the error bar is quietly too
small by $\sqrt{g_{\mathrm{true}}}$. This is not hypothetical; it occurred on the
first attempt to use the machinery on a viscosity run, where $12$ samples per
trajectory returned $g = 1.00$ against a correlation time of $78$ code units and a
sampling interval of $6.7$.

Two guards now fire on that. If a trajectory carries fewer than about $30$
samples, or if the two error estimates disagree by more than a factor of two, the
correction is rejected and the assumption-free between-trajectory error is used
instead. A well-sampled series ($4000$ samples, $\phi = 0.9$) passes both and is
accepted.

\paragraph{Practical consequence.} A production run needs enough \emph{samples} as
well as enough \emph{time}. Sampling a window of $1600$ time units at a
correlation time of $\approx78$ with only $40$ samples gives roughly two samples
per correlation time --- too few to estimate $g$, so the analysis falls back to the
weaker estimator. Around twenty samples per correlation time is a reasonable
target.

% ===========================================================================
\chapter{Verification}
\label{ch:verification}

\section{Unit and regression tests}

Thirteen self-contained executables run under \code{ctest}. Besides the
component tests (RNG stream, config utilities, individual force laws, flow
interpolation, mobility, bending, EV), three deserve mention.

\begin{itemize}
  \item \code{test\_direction} asserts the geometric identity underpinning the
        implicit solve: the updated connector is parallel to $\bm{\Gamma}$.
  \item \code{test\_integrator\_rouse} reproduces the original Fortran
        regression to $10^{-6}$, including two chained integrations on one
        continued RNG stream.
  \item \code{test\_external\_force} checks the applied-force bookkeeping and
        then its effect through the integrator. Measuring $\langle Q_z\rangle$
        directly is noise-limited, but the Hookean dumbbell is \emph{linear}, so
        two runs driven by the same random stream differ deterministically:
        differencing at fixed seed cancels the noise entirely and recovers the
        exact forced extension to six digits ($0.499999$, $1.999997$, $4.999992$
        for $f = 0.5,\,2,\,5$).
  \item \code{test\_hi\_stability} guards the failure mode of
        Sections~\ref{sec:spring-robustness} and~\ref{sec:spectral-bounds}: that
        the implicit solve keeps bonds finitely extensible for
        $\gamma$ up to $10^{20}$; that the Chebyshev square root remains finite
        and satisfies \eqref{eq:fd} on strongly stretched configurations,
        including the collective-translation mode; and that full FENE--Fraenkel
        chains with HI in shear remain bounded over thousands of steps with
        Chebyshev and Cholesky agreeing statistically.
\end{itemize}

\section{Comparison with the original Fortran}
\label{sec:fortran-compare}

The port is validated against the legacy Fortran executable at two levels
(\code{validation/compare\_fortran.py}).

\paragraph{Force kernel --- exact.} The Fortran writes both the configuration and
its own computed total force at each sample. Feeding those configurations through
the C\texttt{++} force routine and comparing reproduces the Fortran forces to
machine precision (worst relative difference $\sim 2\times10^{-16}$ over $10^{3}$
configurations). This check is independent of the random number stream.

\paragraph{Initial distribution --- distributional.} The $t=0$ bond-length
distributions of the two codes are compared by a two-sample
Kolmogorov--Smirnov test, which isolates the initial-configuration sampler from
the dynamics. They are statistically indistinguishable ($p\approx0.19$ over
$2\times10^{4}$ bonds).

\paragraph{Dynamics --- statistical.} The Fortran seeds its generator from the
wall clock, so trajectories cannot be matched one-to-one and only ensemble
averages are meaningful. Over $\sim500$ trajectories per code, the rheological
observables ($\eta_p$, $N_1$) agree within one standard error at every sample
time, and $R^{2}$, $R_g^{2}$ within $0.8$ standard errors. The $x$-stretch, being
an extreme-value statistic, is the noisiest quantity and shows the largest
scatter ($1$--$2\sigma$) without indicating a systematic difference.

One bookkeeping detail: the Fortran time loop is inclusive and overshoots by one
step, so its final sample sits at $T+\dt$ rather than $T$. This is $\sim0.1\%$ of
the integrated time, far below statistical error, and the comparison script
reports it explicitly so that it is not mistaken for a physical discrepancy.

% ===========================================================================
\chapter{Summary of departures from the thesis appendices}
\label{sec:departures}

\begin{longtable}{@{}p{0.28\textwidth}p{0.32\textwidth}p{0.32\textwidth}@{}}
\toprule
Topic & Thesis / legacy Fortran & Current code \\
\midrule
\endhead
Chebyshev spectral bounds
  & Fixman's two-vector Rayleigh-quotient estimate, Eq.~\eqref{eq:fixman}
  & $m$-step Lanczos with full re-orthogonalisation, Sturm bisection for extreme
    Ritz values, outward padding \eqref{eq:padding}, warm-started from the previous
    step's Ritz vectors (Section~\ref{sec:warm-start}). The Fixman estimate is not
    a bound and overestimates $\lambda_{\min}$, which caused divergence for
    stretched chains (Section~\ref{sec:spectral-bounds}). \\
\addlinespace
Square-root failure handling
  & none
  & fluctuation--dissipation check \eqref{eq:fd} with term growth, then exact
    Cholesky fallback for that step. \\
\addlinespace
Cubic root selection
  & closed-form root
  & closed-form root, verified to lie in the physical bracket, with bisection
    fallback if not; result clamped strictly inside the force singularities
    (Section~\ref{sec:spring-robustness}). \\
\addlinespace
Bounded-WLC lookup tables
  & adaptive table built at start-up, \code{hunt}/\code{polint} interpolation
  & not implemented; \code{rtsafe} is called directly each step. \\
\addlinespace
Flow field
  & twelve named flow types plus harmonic traps
  & a single velocity-gradient tensor $\bm{\kappa}$, constant or linearly
    interpolated from a table; traps removed. \\
\addlinespace
Sampling and output
  & inside the integrator, NetCDF
  & owned by the Python driver, which integrates in segments; HDF5. \\


\addlinespace
Bending constant
  & Saadat--Khomami fit only
  & that, plus Langevin, Yamakawa and \code{match\_chain}, which matches the
    chain's $\langle R^2\rangle$ directly (Section~\ref{sec:cg-chain}). \\

\addlinespace
Spring constant
  & unbounded
  & optional ceiling \code{max\_H} with $\sigma$ re-solved
    (Section~\ref{sec:cg-hcap}); $H$ sets the time unit, so an uncapped $H$ makes
    reaching the relaxation time arbitrarily expensive. \\
\addlinespace
Chain friction
  & not addressed here
  & $h^{*}$ solved from the Kirkwood hydrodynamic radius against a measured
    diffusivity, with the full set of unit conversions
    (Chapter~\ref{ch:dynamics}). \\
\addlinespace
Error bars
  & scatter between trajectories
  & autocorrelation-corrected, with blocking and between-trajectory estimates as
    cross-checks and guards against the case where the correction cannot be
    estimated (Chapter~\ref{ch:statistics}). \\
\addlinespace
External forces
  & not present
  & per-bead, constant or time-interpolated (Section~\ref{sec:external}), enabling
    direct force-extension measurement (Chapter~\ref{ch:force-extension}). \\
\addlinespace
Parallelism
  & MPI across ranks
  & multiprocessing over trajectories, seeded by index so results are independent
    of worker count. \\
\bottomrule
\end{longtable}

\paragraph{Unchanged.} The SDE \eqref{eq:sde}, the predictor--corrector scheme
and its $\bm{\Upsilon}$ formulation, all seven spring force laws and their
implicit equations, the RPY tensor, the bending force
\eqref{eq:bending-force}, the excluded-volume forms, the \"{O}ttinger noise
\eqref{eq:ottinger-noise}, the \code{ran\_1} generator, and the
initial-configuration samplers are all mathematically identical to the thesis and
the legacy code, and are verified as such where an oracle exists.

\end{document}
